\documentclass[12pt,a4paper,answers]{exam}

% Packages

%% Language and font encodings
\usepackage[english]{babel}
\usepackage[utf8]{inputenc}
\usepackage[T1]{fontenc}

%% Layout
\usepackage[a4paper,margin=2.5cm]{geometry}
\usepackage{parskip}

%% Fonts
\usepackage{bm}
\usepackage{microtype}
\usepackage{mathrsfs}
\usepackage{lmodern}

%% Core math
\usepackage{amsmath}
\usepackage{amssymb}
\usepackage{amsthm}
\usepackage{mathtools}

%% Useful math
\usepackage{mathdots}
\usepackage{cool}

%% Algorithms and code
\usepackage{algorithm2e}
\usepackage{listings}
\usepackage{lstlang0}

%% Graphics
\usepackage{xcolor}
\usepackage{graphicx}
\usepackage{tikz}
\usetikzlibrary{arrows,calc,decorations.markings,cd}

%% Linking and numbering
\usepackage{enumitem}
\usepackage[hidelinks]{hyperref}

%% References
\usepackage[
backend=bibtex,
style=alphabetic,
sorting=ynt
]{biblatex}

\addbibresource{report.bib}

% Useful commands and definitions

%% Bold question numbers
\renewcommand{\thequestion}{\textbf{\arabic{question}}}

%% Fix spacing between theorems due to parskip

\makeatletter
\def\thm@space@setup{%
  \thm@preskip=\parskip \thm@postskip=0pt
}
\makeatother

%% Listings style

\definecolor{mygreen}{rgb}{0,0.6,0}
\definecolor{mygray}{rgb}{0.5,0.5,0.5}
\definecolor{mymauve}{rgb}{0.58,0,0.82}

\lstset{ %
  backgroundcolor=\color{white},   % choose the background color; you must add \usepackage{color} or \usepackage{xcolor}; should come as last argument
  basicstyle=\footnotesize\ttfamily,            % the size of the fonts that are used for the code
  breakatwhitespace=false,         % sets if automatic breaks should only happen at whitespace
  breaklines=true,                 % sets automatic line breaking
  captionpos=b,                    % sets the caption-position to bottom
  commentstyle=\color{mygreen},    % comment style
  deletekeywords={...},            % if you want to delete keywords from the given language
  escapeinside={\%*}{*)},          % if you want to add LaTeX within your code
  extendedchars=true,              % lets you use non-ASCII characters; for 8-bits encodings only, does not work with UTF-8
  frame=single,	                   % adds a frame around the code
  keepspaces=true,                 % keeps spaces in text, useful for keeping indentation of code (possibly needs columns=flexible)
  keywordstyle=\color{blue},       % keyword style
  language=Octave,                 % the language of the code
  morekeywords={*,...},            % if you want to add more keywords to the set
  numbers=left,                    % where to put the line-numbers; possible values are (none, left, right)
  numbersep=5pt,                   % how far the line-numbers are from the code
  numberstyle=\tiny\color{mygray}, % the style that is used for the line-numbers
  rulecolor=\color{black},         % if not set, the frame-color may be changed on line-breaks within not-black text (e.g. comments (green here))
  showspaces=false,                % show spaces everywhere adding particular underscores; it overrides 'showstringspaces'
  showstringspaces=false,          % underline spaces within strings only
  showtabs=false,                  % show tabs within strings adding particular underscores
  stepnumber=2,                    % the step between two line-numbers. If it's 1, each line will be numbered
  stringstyle=\color{mymauve},     % string literal style
  tabsize=2,	                   % sets default tabsize to 2 spaces
  % title=\lstname                   % show the filename of files included with \lstinputlisting; also try caption instead of title
}

%% Make higher level counter than section
\makeatletter
\newcounter{outersection}
\renewcommand\theoutersection{\@Roman\c@outersection}
\newcommand\outersection{%
   \if@noskipsec \leavevmode \fi
   \par
   \addvspace{4ex}%
   \@afterindentfalse
   \secdef\@outersection\@soutersection}

\def\@outersection[#1]#2{%
  \ifnum \c@secnumdepth >\m@ne
  \refstepcounter{outersection}%
  \addcontentsline{toc}{outersection}{\theoutersection\hspace{1em}#1}%
  \else
  \addcontentsline{toc}{outersection}{#1}%
  \fi
  {\parindent \z@ \raggedright
    \interlinepenalty \@M
    \normalfont
    \ifnum \c@secnumdepth >\m@ne
    \centering\huge\bfseries \outersectionname\ \theoutersection
    \par\nobreak
    \fi
    \LARGE \bfseries #2%
    \markboth{}{}\par}%
  \nobreak
  \vskip 3ex
  \@afterheading}
\def\@soutersection#1{%
  {\parindent \z@ \raggedright
    \interlinepenalty \@M
    \normalfont
    \huge \bfseries #1\par}%
  \nobreak
  \vskip 3ex
  \@afterheading}
\newcommand\outersectionname{Part}
\makeatother

%% The standard sets of numbers
\newcommand{\C}{\mathbb{C}}
\newcommand{\N}{\mathbb{N}}
\newcommand{\Q}{\mathbb{Q}}
\newcommand{\R}{\mathbb{R}}
\newcommand{\Z}{\mathbb{Z}}

%% \abs{} and \norm{} give absolute value and norm respectively
%% \abs*{} and \norm*{} give versions which resize
\DeclarePairedDelimiter\abs{\lvert}{\rvert}
\DeclarePairedDelimiter\norm{\lVert}{\rVert}

\DeclareMathOperator{\Mor}{Mor}
\DeclareMathOperator{\Map}{Map}
\DeclareMathOperator{\im}{im}
\DeclareMathOperator{\Aut}{Aut}
\DeclareMathOperator{\Sym}{Sym}
\DeclareMathOperator{\GL}{GL}
\DeclareMathOperator{\lcm}{lcm}
\makeatletter
\newcommand*{\bdiv}{%
  \nonscript\mskip-\medmuskip\mkern5mu%
  \mathbin{\operator@font div}\penalty900\mkern5mu%
  \nonscript\mskip-\medmuskip
}
\makeatother

\theoremstyle{definition}
\newtheorem{theorem}{Theorem}[section]
\newtheorem{definition}[theorem]{Definition}
\newtheorem{corollary}[theorem]{Corollary}
\newtheorem{lemma}[theorem]{Lemma}

\begin{document}

\begin{titlepage}
  \begin{center}
    \vspace*{1cm}

    \Huge
    \textbf{Computational Group Theory}

    \vspace{2cm}

    \textbf{Peter Huxford}

    \vfill

    A report presented for\\
    COMPSCI 380\\
    Undergraduate Project in Computer Science

    \vspace{0.8cm}

    \includegraphics[width=0.4\textwidth]{UoA.png}

    \Large
    Computer Science\\
    University of Auckland\\
    New Zealand\\
    October 30, 2017

  \end{center}
\end{titlepage}

\outersection*{Introduction}

This report consists of two parts. The first part is a collection of solutions to almost all of the exercises in the first four chapters of the book Presentations of Groups, by D. L. Johnson. This book contains important mathematical foundations for many of the important algorithms in the field of Computational Group Theory. The prerequisite knowledge required to understand these solutions is the contents of the first four chapters of this book, which themselves assume some basic familiarity with group theory. We outline what they discuss below.

The first chapter defines the notion of a free group by its Universal Property, and proceeds to prove some basic properties from this definition. The existence of free groups on a set $X$ is proven by constructing them explicitly, in terms of reduced words. These are strings consisting of characters $x$ and $x^{-1}$ for $x\in X$, which contain no substrings of the form $xx^{-1}$ or $x^{-1}x$. Multiplication and inverses are defined in a natural way to produce a group.

It is illustrated by several examples in the remainder of the chapter, and in the following exercises, that although the Universal Property provides slick proofs of some facts, many results about free groups are most conveniently proven by referring to reduced words.

The second and third chapters provide two substantially different approaches for proving the nontrivial result known as the Nielsen-Schreier Theorem, that every subgroup of a free group is free. The second chapter describes Schreier's method, which provides a fairly explicit description of a basis of such a subgroup in terms of what is known as a Schreier transversal, as well as a formula for the rank of the subgroup in the finite case.

The third chapter discusses Nielsen's method. This approach involves applying a sequence of what are known as Nielsen transformations to a generating set of a subgroup of a free group, to reach a new generating set which is Nielsen-reduced, and hence a basis. Unfortunately, Nielsen's method only proves the Nielsen-Schreier theorem in the case of finitely generated subgroups of a free group. However it provides significant insight into the structure of the automorphism group of the free group.

Finally, the fourth chapter is about group presentations. These are groups defined by some generating set and relations. Various important concepts are introduced, such as von Dyck's theorem and Tietze transformations. Von Dyck's theorem provides a way of proving that a given group presentation is isomorphic to some known group, and Tietze transformations are a formal way of simplifying the relations and generators in a group presentation. The chapter ends with van Kampen diagrams, which are an interesting visual way of demonstrating certain relations in a particular group presentation.

The second part of the report discusses the topic of computation with finitely generated abelian groups. The prerequisite knowledge to understand the main ideas of this part is only introductory group theory, but ideally the reader would have knowledge beyond this.

Restricting oneself to the abelian case does make things simpler, which is helpful because many of the problems one would wish to solve are undecidable in the non-abelian case. For example, there is no algorithm to determine if a given finite presentation has finite order (or even is trivial).

However, such an algorithm is discussed in the second part of this report in the case where the group is also known to be abelian. In fact we give a procedure for which given a description of a finitely generated abelian group $G$, calculates integers $d_1,\ldots,d_r,k$ such that $G\cong\Z_{d_1}\oplus\cdots\oplus\Z_{d_r}\oplus\Z^k$. The description of $G$ is given by an integer matrix, which we transform into a diagonal matrix, known as its Smith Normal Form.

When implementing a naive algorithm inspired by Gaussian elimination, one quickly finds that entries become very large even for relatively small inputs, making calculations in practice far too expensive to carry out. This part of the report discusses some modular techniques in detail, which limit the size of entries, and speed up computation.

\clearpage

\outersection{Solution To Exercises in Presentations of Groups by D. L. Johnson}

%% Sections numbered with Chapter X
\renewcommand{\thesection}{Chapter \arabic{section}}

\section{Free Groups}

\begin{questions}
\question Given a group $G$ and a subset $X\subseteq G$, set $\langle X \rangle = \bigcap\{H\leq G \mid H\supseteq X\}$ and let $W$ be the set of elements of $G$ that can be written as words in $X^{\pm}$. Prove that $W=\langle X \rangle$.
  \begin{solution}
    Let $a,b\in W$, then $a=x_1\cdots x_m$ and $b=y_1\cdots y_n$, for some $x_i\in X^\pm$, $1\leq i\leq m$, and $y_j\in X^\pm$, $1\leq j\leq n$. Then both
    \[ ab=x_1\cdots x_my_1\cdots y_n \quad \text{and} \quad a^{-1}=x_m^{-1}\cdots x_1^{-1} \]
    are words in $X^\pm$, thus $ab\in W$ and $a^{-1}\in W$. Hence $W\leq G$. Moreover each $x\in X$ is a word in $X^\pm$, so $W\supseteq X$. It follows that $\langle X \rangle\subseteq W$.

    Suppose that $H\leq G$ satisfies $H\supseteq X$. We claim that $W\subseteq H$. Let $a\in W$, then $a=x_1\cdots x_m$ for some $x_i\in X^\pm$, $1\leq i\leq m$. For each $1\leq i\leq m$, if $x_i\in X$ then $x_i\in H$. Otherwise $x_i^{-1}\in X$, hence $x_i^{-1}\in H$, and since $H$ is a subgroup $x_i\in H$.

    As $H$ is a subgroup, it follows that the product $a=x_1\cdots x_m$ lies in $H$. Thus $W\subseteq H$, and hence $W\subseteq\langle X \rangle$. The other inclusion yields $W=\langle X \rangle$.
  \end{solution}

\question Let $G$ and $H$ be groups, and $\theta,\phi\colon G\to H$ homomorphisms, and let $X\subseteq G$ be such that $G=\langle X \rangle$. Prove that if $\theta,\phi$ agree on $X$ (i.e. $x\theta=x\phi$, $\forall x\in X$) then $\theta=\phi$ (i.e. $g\theta=g\phi$, $\forall g\in G$).
  \begin{solution}
    For each $x\in X$ we have $x^{-1}\theta = (x\theta)^{-1} = (x\phi)^{-1} = x^{-1}\phi$. Thus $\theta,\phi$ agree on $X^\pm$. Let $a=x_1\cdots x_m$, for $x_i\in X^\pm$, $1\leq i\leq m$, be a word in $X^\pm$. Then
    \[ a\theta = (x_1\cdots x_m)\theta = (x_1\theta)\cdots(x_m\theta) = (x_1\phi)\cdots(x_m\phi) = (x_1\cdots x_m)\phi = a\phi. \]
    Thus $\theta,\phi$ agree on all words in $X^\pm$. So $\theta,\phi$ agree on $\langle X \rangle=G$ by \textbf{1}. Hence $\theta=\phi$.
  \end{solution}

\question Prove directly from the definition that the infinite cyclic group $Z=\{x^n \mid n\in\Z\}$ is free of rank 1.
  \begin{solution}
    We claim $Z$ is free on the subset $\{x\}$ (and hence is free of rank 1). Let $G$ be a group, and let $\theta\colon \{x\}\to G$ be any map. Define $\theta'\colon Z\to G$ by $x^n\theta'=(x\theta)^n$ for $n\in\Z$. Clearly $x\theta'=x\theta$. Additionally $\theta'$ is a homomorphism, as
    \[ (x^mx^n)\theta' = (x^{m+n})\theta' = (x\theta)^{m+n} = (x\theta)^m(x\theta)^n = (x^m\theta')(x^n\theta'). \]
    Every element of $Z$ is a word in $\{x\}^\pm$, so by \textbf{1}, $G=\langle \{x\} \rangle$. Therefore, by \textbf{2} any homomorphism $\phi\colon Z\to G$ with $x\phi=x\theta$ is equal to $\theta'$. So $\theta'$ is the unique such homomorphism, and thus $Z$ is free on $\{x\}$.
  \end{solution}

\question (H. Smith) Let $F$ be free of rank $r$ and $G$ a group isomorphic to $F$. Show that $G$ is free of rank $r$.
  \begin{solution}
    Suppose $F$ is free on the subset $X\subseteq F$, where $\abs{X}=r$, and let $\phi\colon F\to G$ be an isomorphism. We claim that $G$ is free on the subset $Y\coloneqq X\phi\subseteq G$, and hence is free of rank $\abs{Y}=\abs{X}=r$.

    Let $H$ be a group, and let $\theta\colon Y\to H$ be a map. The isomorphism $\phi$ induces a bijection $f\colon X\to Y$. Since $F$ is free on $X$, there is a unique homomorphism $\varphi\colon F\to H$ extending the map $f\theta\colon X\to Y\to H$.

    Let $\theta'\coloneqq\phi^{-1}\varphi\colon G\to F\to H$. Then $\theta'$ is a homomorphism extending $\theta$, because whenever $y\in Y$, $y\phi^{-1}=yf^{-1}\in X$, and thus
    \[ y\theta' = y\phi^{-1}\varphi = (yf^{-1})\varphi = (yf^{-1})f\theta = y\theta.  \]
    Moreover, if $\psi\colon G\to H$ is a homomorphism extending $\theta$, then $\phi\psi\colon F\to G\to H$ extends $f\theta$, because whenever $x\in X$, $x\phi=xf\in Y$, and thus $x\phi\psi = xf\psi = xf\theta$. By the uniqueness of $\varphi$ we have
    \[ \phi\psi = \varphi \implies \psi = \phi^{-1}\varphi = \theta'. \]
    Therefore $\theta'$ is the unique homomorphism extending $\theta$. This shows $G$ is free on $Y$.
  \end{solution}

\question Given a group $G$ and a normal subgroup $N\lhd G$ with $G/N=F$ a free group, prove that $G$ has a subgroup $C$ such that
  \[ N \cap C = \{e\},\ NC = G,\ C\simeq F. \]

  \begin{solution}
    Suppose $F$ is free on the subset $X\subseteq F$, and let $\pi\colon G\twoheadrightarrow G/N=F$ be the canonical projection. Note that the elements of $F$ are cosets of $N$ in $G$. By the Axiom of Choice there is a function $\theta\colon X\to G$ with $x\theta\in x$ for each $x\in X$. This implies $x\theta\pi=x$ for each $x\in X$.

    Since $F$ is free on $X$, $\theta$ extends to a homomorphism $\theta'\colon F\to G$. If $x\in X$, then
    \[ x^{-1}\theta'\pi = (x\theta)^{-1}\pi = (x\theta\pi)^{-1} = x^{-1}. \]
    Therefore $x^{-1}\theta'\in x^{-1}$ for each $x\in X$, so $x\theta'\in x$ and $x\theta'\pi=x$ for each $x\in X^\pm$. Now if $x_1,\ldots,x_m\in X^\pm$, then
    \[ (x_1\cdots x_m)\theta'\pi = (x_1\theta')\cdots(x_m\theta')\pi = (x_1\theta'\pi)\cdots(x_m\theta'\pi) = x_1\cdots x_m. \]
    Hence $x\theta'\pi=x$ for all words $x$ in $X^\pm$. Since $F$ is free on $X$, $F=\langle X \rangle$, and so by \textbf{1}, $x\theta'\pi=x$, for all $x\in F$. Let $Y\coloneqq X\theta'$, and $C\coloneqq F\theta'\leq G$. We prove that $\pi{\restriction_C}$ and $\theta'$ are two-sided inverses of each other.

    We know $\theta'\pi$ is the identity on $F$. If $c\in C$ then $c=x\theta'$ for some $x\in F$. Therefore $c\pi\theta'=x\theta'\pi\theta'=x\theta'=c$, so $\pi{\restriction_C}\theta'$ is the identity on $C$. Hence $\pi{\restriction_C}$ is an isomorphism from $C$ to $F$, and so $C\simeq F$. Moreover,
    \[ N\cap C = \ker\pi\cap C = \ker\pi{\restriction_C}=\{e\}. \]
    Let $g\in G$, and set $c=g\pi\theta'$. Then $c\in C$, and
    \[ gc^{-1}\pi=(g\pi)(c\pi)^{-1}=(g\pi)(g\pi\theta'\pi)^{-1}=(g\pi)(g\pi)^{-1}=e. \]
    So $gc^{-1}\in\ker\pi=N$, hence $g\in NC$ and thus $NC=G$. This completes the proof.
  \end{solution}

\question How many reduced words are there of length $l$ in the free group of rank $r$? How many of these are cyclically reduced?
  \begin{solution}
    There are $2r$ possibilities for the first character, and $2r-1$ possibilities for each subsequent character. Hence there are $2r(2r-1)^{l-1}$ such words.
  \end{solution}

\question Show that the free abelian group of rank $r$ is isomorphic to the direct product of $r$ copies of $Z$.

\question Let $F$ be free of rank $r$. Show that $F/F'$ is free abelian of rank $r$.

\question Prove that the free group $F$ of finite rank $r$ has only a finite number of normal subgroups of index $n\in\N$. [\emph{Hint:} Consider homomorphisms from $F$ onto groups $G$ of order $n$.] Is the same true for arbitrary subgroups (not necessarily normal) of index $n$?
  \begin{solution}
    A given normal subgroup $N\lhd F$ of index $n$ induces a homomorphism $F\twoheadrightarrow F/N$ with kernel $N$. Since $F/N$ is a group of order $n$, the (faithful) regular action of $F/N$ on itself yields an embedding $F/N\hookrightarrow S_n$. Composing these maps produces a homomorphism $\phi\colon F\to S_n$ with $\ker\phi=N$.

    Therefore it suffices to prove that $\Mor(F,S_n)$ is finite. Suppose $F$ is free on $X$, where $\abs{X}=r$. Then $\abs{\Mor(F,S_n)}=\abs{\Map(X,S_n)}=\abs{X}^{\abs{S_n}}=r^{n!}$ is finite.
  \end{solution}

\question Let $F$ be free on a subset $X\subseteq F$, and let $Y$ be a subset of $X$. Prove that the subgroup $H=\langle Y \rangle\leq F$ is free on the set $Y$.

\question Let $F$ be free on $\{x,y\}$ and define $a_n=x^{-n}yx^n$, $n\in\N$. Prove that no reduced word in $a_1,\ldots,a_r$ ($r\in\N$) can be equal to $e$, and deduce that the group $F_r\coloneqq\langle a_1,\ldots,a_r \rangle$ is free on these generators. Show that the group $F_0=\bigcup_{r\in\N} F_r = \langle \{x_n \mid n\in\N\} \rangle$ is free of rank $\aleph_0$.
  \begin{solution}
    A reduced word of positive length in $\{a_1,\ldots,a_r\}^\pm$ can be written in the form $a_{i_1}^{e_1}\cdots a_{i_m}^{e_m}$ for $m>0$, $e_j\in\Z\setminus\{0\}$, $i_j\in\{1,\ldots,r\}$, with $i_j\neq i_{j+1}$, by grouping adjacent terms into powers appropriately.

    It is easily verified that $a_n^e=x^{-n}y^ex^n$, hence
    \begin{align*}
      a_{i_1}^{e_1} \cdots a_{i_m}^{e_m} &= (x^{-i_1}y^{e_1}x^{i_1})\cdots(x^{-i_m}y^{e_m}x^{i_m}) \\
                                         &= x^{-i_1}y^{e_1}x^{i_1-i_2}y^{e_2}\cdots x^{i_{m-1}-i_m}y^{e_m}x^{i_m}.
    \end{align*}
    This is an alternating product of nonzero powers of $x,y$ with $2m+1>0$ terms. Therefore as a reduced word in $\{x,y\}^\pm$ it is of positive length, and so nontrivial. Since $a_1,\ldots,a_r$ generate $F_r$, it is free on these generators.

    In fact, any reduced word of positive length in $\{a_n \mid n\in\N\}^\pm$ is a reduced word in $\{a_1,\ldots,a_r\}^\pm$ for some $r\in\N$ (because it has only finitely many terms), and is thus non-trivial. Thus $F_0$ is free on the generators $\{a_n\mid n\in\N\}$, and is therefore free of rank $\abs{\{a_n\mid n\in\N\}}=\aleph_0$.
  \end{solution}

\question Prove that two elements $a,b$ in a free group $F$ are conjugate in $F$ if and only if $\check{a}$ and $\check{b}$ (as defined in (5)) are cyclic rearrangements of one another, that is, if $\check{a}=x_1\ldots x_l$, then $\check{b}=x_r\ldots x_lx_1\ldots x_{r-1}$ for some $r$, $1\leq r\leq l$.

\question Prove that, in a free group $F$, no non-trivial element can be conjugate to its inverse. Deduce that if $H=\langle a \rangle$ is an infinite cyclic subgroup of $F$, then $N(H)=C(a)$, where $N(H)=\{w\in F \mid w^{-1}Hw=H \}$ is the normalizer of $H$ in $F$.

\question Show that no word $w\neq e$ in a free group can be conjugate to a proper power of itself, namely, to $w^n$ where $n\in\Z$ and $n\neq1$.

\question Prove that the set of reduced words of even length in $F=F(\{x,y\})$ forms a subgroup, call it $H$. Show that $H$ is generated by the elements
  \[ a = x^2, b = xy, c = xy^{-1}. \]
  Verify that if $w$ is one of the 30 reduced words of length 2 in $\{a,b,c\}^\pm$, then the process of expressing $w$ as a reduced word in $\{x,y\}^\pm$ cannot involve the cancellation of the second letter of $a$, $b$ or $c$, nor of the first letter of $a^{-1}$, $b^{-1}$, $c^{-1}$. Thus show that if $w$ is a reduced word of length $l$ in $\{a,b,c\}^{\pm}$, then as a reduced word in $\{x,y\}^\pm$, its length is at least $l$. Deduce that $H$ is a free group of rank 3.

\question Generalise Exercise 15 from rank 2 to rank $r$, that is, find a set of free generators for the group of reduced words of even length in $F(\{x_1,\ldots,x_r\})$.

\question For a fixed $l,r\in\N$, what can you say about the group generated by the reduced words of length $l$ in the free group of rank $r$?

\question The \emph{centre} of any group $G$ is defined by
  \[ Z(G) \coloneqq \{ z\in G \mid zg=gz,\ \forall g\in G\}. \]
  Prove that the center of every non-cyclic free group is trivial.

\question Prove that no finite group can be the union of the conjugates of a proper subgroup.

\question (J. Wiegold) In the free group $F=F(\{a,b\})$ of rank 2, let
  \[ A = \{ a_k \mid k\in\N \} \]
  be an enumeration of those reduced words that begin and end with a power of $a$. Prove that if
  \[ B = \{b^{-k}a_kb^k \mid k\in\N \}, \]
  then the subgroup $H=\langle B \rangle$ is free on $B$. Deduce that $H<F$, but $\bigcup_{w\in F} H^w=F$.
\end{questions}

%%% Local Variables:
%%% mode: latex
%%% TeX-master: "johnson"
%%% End:
\section{Schreier's Method}
\begin{questions}
\question Prove that the relation $<$ of Definition 2 satisfies (o1), (o2), (o4).
  \begin{solution}
    Let $s=x_1\cdots x_l$, $t=y_1\cdots y_m$, $u=z_1\cdots z_n$ be reduced words in $X^\pm$. Since $l\not <l$, and $x_i=x_i$ for $1\leq i\leq r$, we have not ($s<s$), thus (O1) is satisfied.

    Suppose that $s<t$ and $t<u$. Then $l\leq m\leq n$, so if either $l<m$ or $m<n$, then $l<n$, and so $s<u$. Otherwise $l=m=n$, $x_r<y_r$ for $r=\min\{i \mid x_i\neq y_i\}$, and $y_{r'}<z_{r'}$ for $r'=\min\{j \mid y_j\neq z_j\}$. Note $x_i=y_i=z_i$ for all $1\leq i<\min\{r,r'\}$

    If $\min\{r,r'\}=r$, then $x_r<y_r\leq z_r$, so that $s<u$. If $\min\{r,r'\}=r'$, then $x_{r'}\leq y_{r'}<z_{r'}$, so that $s<u$. In all cases $s<u$, so (O2) is satisfied.

    Let $S$ be a non-empty subset of $F$, and define $n=\min\{l(s) \mid s\in S\}\subseteq\N$. We inductively define a sequence $S_1,\ldots,S_n$ of non-empty subsets of $X^\pm$. Let $i\geq1$, and suppose $S_1,\ldots,S_{i-1}$ have been defined. Set $x_j=\min S_j$, for $1\leq j<i$. Define
    \[ S_i = \{ x\in X^\pm \mid x_1\cdots x_{i-1}xy_{i+1}\cdots y_n \in S\text{ is reduced, for some } y_j\in X^\pm \}. \]
    Since $S$ contains a reduced word of length $n$, it follows that $S_1$ is non-empty. For $i>1$, suppose that $S_j$ is non-empty for each $1\leq j<i$. Then the definitions $x_j=\min S_j$ make sense, and moreover for each $x\in S_{i-1}$, there exists a reduced word $x_1\cdots x_{i-2}xy_i\cdots y_n\in S$. Taking $x=x_{i-1}$ shows that $S_i$ is non-empty.

    If we finally define $x_n=\min S_n$, then we claim that the word $s=x_1\cdots x_n$ is the least element of $S$. By definition, $x_1\cdots x_{n-1}x$ is reduced and in $S$ for each $x\in S_n$. Taking $x=x_n$ shows that $s$ is a reduced word of length $n$. Let $t=y_1\cdots y_m\in S$ be a reduced word. Then $n\leq m$ by definition of $n$. If $n<m$, then $s<t$.

    Suppose that $n=m$, and that $s\neq t$. Let $r=\min\{i \mid x_i\neq y_i\}$. Then $x_r\neq y_r\in S_r$, and thus $x_r=\min S_r<y_r$, hence $s<t$. It follows that $S$ has a least element, $s$. Therefore (O4) is satisfied.
  \end{solution}

\question What is the 666th element of $F=F(\{x,y\})$ in the ordering (1)?
  \begin{solution}
    By Exercise \textbf{1.6}, for $l\geq1$ there are $4\cdot3^{l-1}$ words of length $l$ in $F$, and hence there are
    \[ 1+4(1+\cdots+3^{l-1})=1+4\cdot\frac{3^l-1}{3-1}=2\cdot3^l-1 \]
    words of length at most $l$ in $F$. Thus there are $2\cdot3^5-1=485$ words of length at most 5 in $F$. The ordering the question refers orders the elements first by length, and then lexicographically according to the ordering $x<y<x^{-1}<y^{-1}$ on $X^\pm$.

    There are $2\cdot 3^4=162$ reduced words of length 6 beginning with either $x^2$ or $xy$, as in each case there are 3 choices for each of the 4 subsequent letters. Similarly there are $2\cdot3^2=18$ reduced words of length 6 beginning with either $xy^{-1}x^2$ or $xy^{-1}xy$.

    The words described so far make up the first $485+162+18=665$ words in $F$. Therefore the 666th word in $F$ is $xy^{-1}xy^{-1}x^2$.
  \end{solution}

\question With the same $F$ and $<$, what is the second smallest element of the commutator subgroup $F'$?
  \begin{solution}
    By Exercise \textbf{1.8}, $F/F'$ is free abelian of rank 2 on $\{F'x,F'y\}$, and by Exercise \textbf{1.7}, $F/F'\cong Z\times Z$, where $Z$ is the infinite cyclic group $\{x^n \mid n\in\Z\}$, via the isomorphism determined by $F'x\mapsto(x,1)$ and $F'y\mapsto(1,x)$.

    The composite mapping $F\twoheadrightarrow F/F'\to Z\times Z$ has kernel $F'$, and is determined by $x\mapsto(x,1)$ and $y\mapsto(1,x)$. It follows that this map sends a word $w$ to $(x^{e_1},x^{e_2})$, where $e_1$ and $e_2$ are the sums of exponents of $x$-terms and $y$-terms in the word $w$.

    Therefore, elements of $F$ are reduced words which have the same number of occurrences of $x$ and $x^{-1}$ letters, and the same number of $y$ and $y^{-1}$ letters. Thus, the second smallest element of $F'$ is $xyx^{-1}y^{-1}$.
  \end{solution}

\question With the same $F$ and $<$ again, let $H=\overline{S}$ with $S=\{x^4, y^2, (xy)^2\}$. Show that $\abs{F:H}=8$ (using $D_4$) and write down a Schreier transversal for $H$ in $F$. Compile a table like Table 1, and thus obtain a basis $B$ of $H$.
  \begin{solution}
    We think of $D_4=\langle r\coloneqq(1234), s\coloneqq(12)(34) \rangle$ as a subgroup of $S_4$. The map $\theta\colon\{x,y\}\to D_4$ given by $\theta x=r$ and $\theta y=s$ extends to a homomorphism $\theta'\colon F\to D_4$. Since $r,s\in\im\theta'$, and $D_4=\langle r,s \rangle$, it follows that $\im\theta'=D_4$.

    It is simply verified that $S\subseteq\ker\theta'$, and since $\ker\theta'\lhd F$ it follows that $H=\overline{S}\leq\ker\theta'$. Therefore, applying the First Isomorphism Theorem
    \[ \abs{F:H} = \abs{F:\ker\theta'}\cdot\abs{\ker\theta':H} \geq \abs{F:\ker\theta'} = \abs{\im\theta'}= \abs{D_4} = 8. \]

    On the other hand, the elements $x'=Hx$, $y'=Hy$ generate $F/H$, have orders 4 and 2 respectively, and satisfy the relation $(x'y')^2=e\implies y'x'=x'^3y'$. Using this second relation we can write any element of $F/H$ as a word $x'^iy'^j$, and by using the orders of $x'$ and $y'$ we can ensure that $i\in\{0,1,2,3\}$ and $j\in\{0,1\}$.

    Thus, any element of $F/H$ is one of the eight cosets given by $Hx^iy^j$, where $i\in\{0,1,2,3\}$ and $y\in\{0,1\}$. Since $\abs{F:H}\geq8$, these eight cosets must be distinct, and thus $T=\{e,x,x^2,x^3,y,xy,x^2y,x^3y\}$ is a transversal for $H$ in $F$.

    Although $T$ has the Schreier property, taking the least element of each coset should lead to simpler computations for a basis $B$ of $H$. Using the relations between $x'$ and $y'$, we see that $Hx^3=Hx^{-1}$ and $Hx^3y=Hyx$. Therefore $U=\{e,x,y,x^{-1},x^2,xy,yx,x^2y\}$ is a transversal for $H$ in $F$. In fact it has the Schreier property, and thus is a Schreier transversal for $H$ in $F$.

    The table below has rows indexed by $u\in U$ and columns indexed by $z\in\{x,y\}^\pm$, and contains the generator $uz\overline{uz}^{-1}$ in the $(u,z)$ place. The identification of the element $\overline{uz}$ can be done by using the relations between $x',y'$ described earlier, and the other computations are mechanical, as they are just manipulations of words in the free group.

    \begin{center}
      \begin{tabular}{|l||p{2cm}|p{2cm}|p{2.5cm}|p{2.5cm}|}
        \hline
        & $x$ & $y$ & $x^{-1}$ & $y^{-1}$ \\
        \hline
        \hline
        $e$ & $e$ & $e$ & $e$ & $y^{-2}$ \\
        \hline
        $x$ & $e$ & $e$ & $e$ & $xy^{-2}x^{-1}$ \\
        \hline
        $y$ & $e$ & $y^2$ & $yx^{-1}y^{-1}x^{-1}$ & $e$ \\
        \hline
        $x^{-1}$ & $e$ & $x^{-1}yx^{-1}y^{-1}$ & $x^{-4}$ & $x^{-1}y^{-1}x^{-1}y^{-1}$ \\
        \hline
        $x^2$ & $x^4$ & $e$ & $e$ & $x^2y^{-2}x^{-2}$ \\
        \hline
        $xy$ & $xyxy^{-1}$ & $xy^2x^{-1}$ & $xyx^{-1}y^{-1}x^{-2}$ & $e$ \\
        \hline
        $yx$ & $yx^2y^{-1}x^{-2}$ & $yxyx$ & $e$ & $yxy^{-1}x$ \\
        \hline
        $x^2y$ & $x^2yxy^{-1}x^{-1}$ & $x^2y^2x^{-2}$ & $x^2yx^{-2}y^{-1}$ & $e$ \\
        \hline
      \end{tabular}
    \end{center}

    The table entries generate $H$, and the non-identity elements in the first two columns
    \[ B = \{x^4, y^2, yxyx, xyxy^{-1}, xy^2x^{-1}, x^2y^2x^{-2}, x^{-1}yx^{-1}y^{-1}, yx^2y^{-1}x^{-2}, x^2yxy^{-1}x^{-1} \} \]
    form a free basis of $H$.
  \end{solution}

\question Prove that the permutations $x_1=(12)$, $x_2=(23)$, $x_3=(34)$ generate the symmetric group $S_4$. Use the transversal $U=\{e,x_1\}$ of the alternating group $A_4$ to find Schreier generators for $A_4$.
  \begin{solution}
    Note that $x_1x_2x_1=(13)$, $(13)x_3(13)=(14)$, and $x_2x_3x_2=(24)$, hence $\langle x_1,x_2,x_3 \rangle$ contains all transpositions. Since the transpositions generate the symmetric group, it follows that $x_1,x_2,x_3$  generate $S_4$.

    If we set $X=\{x_1,x_2,x_3\}$, then the Schreier generators for $A_4$ are given by
    \begin{align*}
      B &= \{ ux\overline{ux}^{-1} \mid u\in U, x\in X, ux\notin U \} \\
        &= \{ x_2x_1^{-1}, x_3x_1^{-1}, x_1x_2, x_1x_3 \} \\
        &= \{ (123), (12)(34), (132) \}.
    \end{align*}
  \end{solution}

\question Convince yourself that $U=\{x^ky^l \mid k,l\in\Z\}$ is a Schreier transversal for $F'$ in $F=F(\{x,y\})$. Write down the corresponding Schreier generators for $F'$ and express the commutator $x^{-2}y^{-3}x^2y^3$ in terms of them.
  \begin{solution}
    It is clear $U$ has the Schreier property. According to Exercise \textbf{1.8} $F/F'$ is free abelian on $x'=F'x$ and $y'=F'y$, and is isomorphic to $Z\times Z$ via $x'\mapsto (x,1)$ and $y'\mapsto(1,x)$, where $Z=\{x^n \mid n\in\Z\}$ is the infinite cyclic group.

    Thus, every element of $F/F'$ can be written uniquely in the form $(x')^k(y')^l=F'x^ky^l$, for some $k,l\in\Z$. It follows that $U$ is a Schreier transversal for $F'$ in $F$. Let $u\in U$, then $u=x^ky^l$ for some $k,l\in\Z$. Note $uy=x^ky^{k+1}\in U$, however $ux=x^ky^lx\notin U$.

    Since $F/F'$ is abelian, $F'ux=F'xu$. Since $xu=x^{k+1}y^l\in U$, we have $\overline{ux}=xu$. So the set $B$ of Schreier generators for $F'$ is given by
    \[ B = \{ ux\overline{ux}^{-1} \mid u\in U \} = \{ uxu^{-1}x^{-1} \mid u\in U \} = \{ x^ky^lxy^{-l}x^{-k-1} \mid k,l\in\Z \}. \]
    Let $u_1=e$, and inductively define $u_{i+1}=\overline{u_ix_i}$, where $x_i$ is the $i$th letter in $w=x^{-2}y^{-3}x^2y^3$. Then $u_1,\ldots,u_6$ are the truncations $e,x^{-1},\ldots,x^{-2}y^{-3}$ of $w$. We have $u_7=\overline{u_6x}=x^{-1}y^{-3}$, $u_8=\overline{u_7x}=y^{-3}$, and then finally $u_9,u_{10},u_{11}$ are $y^{-2},y^{-1},e$, respectively.

    Following the proof that $B$ generates $F'$, we find the below expression of $[x^2,y^3]$
    \[ x^{-2}y^{-3}x^2y^3 = (x^{-2}y^{-3}xy^3x)(x^{-1}y^{-3}xy^3). \]
  \end{solution}

\question Prove that every subgroup of finite index in a finitely-generated group is finitely generated.
  \begin{solution}
    Let $G$ be generated by a finite set $X$, and suppose $H\leq G$ has finite index. If $U$ is a transversal for $H$ in $G$, then $U$ is finite. If $e\in U$, then the set
    \[ B=\{ux\overline{ux}^{-1} \mid u\in U, x\in X, ux\notin U\} \]
    generates $H$. Clearly $\abs{B}\leq\abs{U}\abs{X}$, so $B$ is finite. Hence $H$ is finitely generated.
  \end{solution}

\question Express each of the seven free generators in the left half of Table 1 as products of conjugates of $a=x^3$, $b=y^2$, $c=x^{-1}y^{-1}xy$.
  \begin{solution}
    The seven free generators are $x^3$, $x^{-1}yxy^{-1}$, $yxy^{-1}x^{-1}$, $y^2$, $xyx^2y^{-1}$, $xy^2x^{-1}$, and $yx^{-1}yx$. We have
    \begin{align*}
      x^3 &= a \\
      x^{-1}yxy^{-1} &= yc^{-1}y^{-1} \\
      yxy^{-1}x^{-1} &= (yx)c^{-1}(yx)^{-1} \\
      y^2 &= b \\
      xyx^2y^{-1} &= (xy)ac(xy)^{-1} \\
      xy^2x^{-1} &= xbx^{-1} \\
      yx^{-1}yx &= (ycy^{-1})(x^{-1}bx).
    \end{align*}
  \end{solution}

\question Prove that the free group of rank 2 contains a free subgroup of any given finite rank as a normal subgroup.
  \begin{solution}
    Let $F=F(\{x,y\})$, and let $C_n=\langle x \rangle$ be the cyclic group of order $n$. Define $\theta\colon\{x,y\}\to C_n$ by $x\theta=x$, $y\theta=1$. Let $\theta'\colon F\to C_n$ be the extension of $\theta$. Then $\ker\theta'\lhd F$, and since $x\in\im\theta'$, $\theta'$ is surjective. Thus $F/\ker\theta'\cong C_n$, and hence $\ker\theta'$ has index $\abs{C_n}=n$.

    From the Nielsen-Schreier Theorem, $\ker\theta'$ is free of rank $(2-1)n+1=n+1$. Thus $F$ contains normal subgroups which are free of any given finite rank $\geq2$.
  \end{solution}

\question Prove Euler's formula for planar graphs: let $\Gamma$ have $v$ vertices, $e$ edges (all straight line segments) and $f$ faces (polygonal regions bounded by edges of $\Gamma$ and containing no vertices of $\Gamma$ in their interiors). Then $v-e+f$ is equal to the number of connected components of $\Gamma$.
  \begin{solution}
    We prove the result by induction on the number of edges of $\Gamma$. If $e=0$, then $v$ is the number of connected components of $\Gamma$, and $f=0$, thus the formula holds. Suppose $e\geq1$, and that the result holds for all graphs with fewer edges.

    Suppose $\Gamma$ has $k$ connected components. Let $\Gamma'$ be the graph $\Gamma$ but with one edge removed, and suppose $\Gamma'$ has $v'$ vertices, $e'$ edges and $f'$ faces, and $k'$ connected components. Then $e'=e-1$ and $v'=v$. By hypothesis $v'-e'+f'=k'$.

    If removing the edge disconnects a connected component of $\Gamma$, then $k'=k+1$, and $f'=f$. Thus $v-(e-1)+f=k+1$, and hence $v-e+f=k$.

    Otherwise, removing the edge doesn't disconnect a component of $\Gamma$, so that $k'=k$. If the removed edge borders a face and the exterior of $\Gamma$, then that face becomes part of the exterior in $\Gamma'$. Otherwise it separates two faces, which combine into a single face in $\Gamma'$. In each case, $f'=f-1$. Hence $v-e+f=v-(e-1)+(f-1)=k$.

    It follows by induction that Euler's formula holds in all planar graphs.
  \end{solution}

\question Use the Nielsen-Schreier theorem and the result of Exercise 1.18 to give a quick proof of Theorem 1.2.
  \begin{solution}
    Let $F$ be a free group and let $w\in F\setminus\{e\}$. By the Nielsen-Schreier theorem, $C(w)\leq F$ is a free group. Clearly $w\in Z(C(w))$, thus by Exercise 1.18, $C(w)$ is a cyclic group. It must be infinite, because $e\neq w\in C(w)$, and no non-trivial power of $w$ is $e$, by the freeness of $F$.
  \end{solution}

\question Let $F=F(X)$ be an arbitrary free group, and $H\leq F$ the subgroup consisting of words of even length. Use Schreier's method to obtain a basis for $H$ (cf. Exx. 1.15, 1.16).
  \begin{solution}
    Let $x'\in X$, then $U=\{e,x'\}$ is a Schreier transversal for $H$ in $F$. For each $x\in X$, we have $\overline{ex}=x'$ and $\overline{x'x}=e$. Therefore $ex\in U$ iff $x=x'$, and $x'x\in U$ iff $x=x'^{-1}$. It follows that
    \begin{align*}
      B &= \{ ux\overline{ux}^{-1} \mid u\in U, x\in X, ux\notin U \} \\
        &= \{ xx'^{-1} \mid x\in X\setminus\{x'\} \} \cup \{ x'x \mid x\in X\setminus\{x'^{-1}\} \},
    \end{align*}
    is a basis for $H$ in $F$. By taking the inverse of some of these elements we obtain a simpler basis $\{x'x \mid x\in X^\pm\setminus\{x'^{-1}\}\}$.
  \end{solution}

\question (M. Hall, Jr.) Let $F=F(X)$ be a free group of finite rank $\abs{X}=r$, and let $n\in\N$ be a fixed natural number. Prove that $F$ contains only finitely many $n$-element Schreier sets, and that if $U$ is one of these, the number of functions $U\times X\to U$ is finite. Deduce that $F$ has only finitely many subgroups of index $n$.
  \begin{solution}
    If $U$ is a Schreier set with $n$ elements, then $U$ contains no words of length $n$ or greater. For, if $x_1\cdots x_n\in U$, then $x_1\cdots x_j\in U$ for each $0\leq j\leq n$. By Exercise 1.6, the set $S$ of words in $X^\pm$ of length $<n$ is finite. The number of $n$-element subsets of $S$ is at most $2^{\abs{S}}$, and so is also finite.

    Hence $F$ contains only finitely many $n$-element Schreier sets $U$. If $U$ is one of these, then since $U$ and $X$ are finite, there are finitely many functions $U\times X\to U$. Thus there are finitely many pairs $(U,f)$ such that $\abs{U}=n$, $f\colon U\times X\to U$, and $U$ is a Schreier transversal of $H\coloneqq\langle \{ux (f(u,x))^{-1} \mid u\in U, x\in X\} \rangle$, and $f(u,x)=\overline{ux}$.

    Let $A$ be the set of all such pairs, and let $B$ be the set of all index $n$ subgroups of $G$. Then the map $A\to B$ given by $(U,f)\mapsto \langle \{ux(f(u,x))^{-1} \mid u\in U, x\in X\} \rangle$ is surjective. For, if $H\leq F$ has index $n$ then $H$ has a Schreier transversal $U$ of size $n$, and the map $(u,x)\mapsto\overline{ux}$ satisfies $H=\langle \{ux\overline{ux}^{-1} \mid u\in U, x\in X\} \rangle$.

    Since $A$ is finite, and we have a surjective map $A\to B$, it follows that $\abs{B}$ is finite. Hence $F$ has finitely many subgroups of index $n$.
  \end{solution}

\question Use the previous exercise to show that a finitely-generated group has only finitely many subgroups of a given finite index.
  \begin{solution}
    Suppose that $G=\langle X \rangle$ is generated by a finite set $X$. Define $F\coloneqq F(X)$, and let $\pi\colon F\to G$ be the extension of the inclusion map $\theta\colon X\to G$. Suppose $H\leq G$ has finite index $n$ in $G$. $G$ acts by right multiplication on the set of cosets $G/H$. Since $\abs{G/H}=n$, this induces a homomorphism $\varphi\colon G\to S_n$.

    Let $N\coloneqq\ker\varphi$. Let $n\in N$, then $Hn=H$, and thus $n\in H$. So $N\leq H$. Let $N',H'\leq F$ be the unique subgroups such that $N'\pi=N$ and $H'\pi=H$. Then $N'\leq H'$, and in fact $N'=\ker(\pi\varphi)$. Thus $F/N'$ is isomorphic to a subgroup of $S_n$.

    Therefore $N'$ has index $\leq n!$ in $F$, and thus $H'$ has index $\leq n!$ in $F$. By \textbf{13}, $F$ contains finitely many subgroups of index $\leq n!$. Therefore $G$ contains finitely many subgroups of index $n$.
  \end{solution}

\end{questions}

%%% Local Variables:
%%% mode: latex
%%% TeX-master: "johnson"
%%% End:
\section{Nielsen's Method}
\begin{questions}
\question By expressing the transposition: $(u_1,\ldots,u_i,\ldots,u_j,\ldots,u_n)\mapsto(u_1,\ldots,u_j,\ldots,u_i,\ldots,u_n)$ as a product of six elementary Nielsen transformations, prove that any permutation of $U$ can be effected by a regular Nielsen transformation.
  \begin{solution}
    We have
    \begin{align*}
      & \phantom{{}={}} (u_1,\ldots,u_i,\ldots,u_j,\ldots,u_n)(ij)(i')(ji)(j')(ij)(i') \\
      &= (u_1,\ldots,u_iu_j,\ldots,u_j,\ldots,u_n)(i')(ji)(j')(ij)(i') \\
      &= (u_1,\ldots,u_j^{-1}u_i^{-1},\ldots,u_j,\ldots,u_n)(ji)(j')(ij)(i') \\
      &= (u_1,\ldots,u_j^{-1}u_i^{-1},\ldots,u_i^{-1},\ldots,u_n)(j')(ij)(i') \\
      &= (u_1,\ldots,u_j^{-1}u_i^{-1},\ldots,u_i,\ldots,u_n)(ij)(i') \\
      &= (u_1,\ldots,u_j^{-1},\ldots,u_i,\ldots,u_n)(i') \\
      &= (u_1,\ldots,u_j,\ldots,u_i,\ldots,u_n).
    \end{align*}
    Thus the transposition $(u_1,\ldots,u_i,\ldots,u_j,\ldots,u_n)\mapsto(u_1,\ldots,u_j,\ldots,u_i,\ldots,u_n)$ can be written as the product $(ij)(i')(ji)(j')(ij)(i')$ of six elementary Nielsen transformations. Since every permutation can be written as a product of transpositions, it follows that any permutation of $U$ can be effected by a regular Nielsen transformation.
  \end{solution}

\question Let $\tau$ be the transformation of $U$ sending $u_i$ to $(u_i^\gamma u_j^\delta)^\varepsilon$, where $i\neq j$ and $\gamma,\delta,\varepsilon\in\{\pm1\}$, and fixing all other $u_k$. Show that $\tau$ is a regular Nielsen transformation.
  \begin{solution}
    Let $\tau_1,\pi_1$ be the identity Nielsen transformations, and $\sigma_1=(ij)$. Define $\tau_{-1}=(i')$, $\sigma_{-1}=(j')(ij)(j')$, and $\pi_{-1}=(i')$. Then $\tau_\gamma,\sigma_\delta,\pi_\varepsilon$ fix all $u_k$ with $k\neq i$. Moreover $\tau_\gamma$ sends $u_i$ to $u_i^{\gamma}$, $\sigma_\delta$ sends $u_i$ to $u_iu_j^{\delta}$, and $\pi_\varepsilon$ sends $u_i$ to $u_i^{\varepsilon}$. Therefore $\tau=\tau_\gamma\sigma_\delta\pi_\varepsilon$ is a regular Nielsen transformation.
  \end{solution}

\question Prove that, except in the trivial case $U=\{e\}$, (N3)$\implies$(N1).
  \begin{solution}
    If (N1) is false, then $e\in U$. Suppose then that there is an $a\in U$, $a\neq e$. Then $l(aea^{-1})=l(e)=0$, but $l(a)-l(e)+l(a^{-1})=2l(a)>0$, and $a,e,a^{-1}\in U^\pm$ with $a\neq e\neq a^{-1}$. Thus (N3) is false. Taking the contrapositive gives the result.
  \end{solution}

\question Let $U=(u_1,\ldots,u_n)$ be a basis for a subgroup $H$ of $F$, and $\tau$ a regular Nielsen transformation of $U$. Prove that $U\tau$ is a basis for $H$. (Hint: for $w\in H$, let $l(w)$, $l_\tau(w)$ denote the length of $w$ as a reduced word in $U$, $U\tau$ respectively. Then prove that $l_\tau(w)\geq l(w)/k$, when $\tau$ is an elementary Nielsen transformation of type (Tk), $k=1,2$.)

\question Prove that the ordering $\ll$ in the proof of Theorem 1 is a well-ordering on the pairs $w^\pm$ in $F\setminus E$.
  \begin{solution}
    The ordering is defined as $w_1^\pm \ll w_2^\pm$ if and only if
    \begin{enumerate}[label=(\alph*)]
    \item $l(w_1)<l(w_2)$, or these are equal and either
    \item $\min\{L(w_1),L(w_1^{-1})\}<\min\{L(w_2),L(w_2^{-1})\}$, or
    \item these mins are equal, and $\max\{L(w_1), L(w_1^{-1})\}<\max\{L(w_2), L(w_2^{-1})\}$,
    \end{enumerate}
    where $L(w)$ is the first $\lfloor (l(w)+1)/2 \rfloor$ letters of $w$. Since $l(w)=l(w^{-1})$, (a) is well-defined on the pairs $w^\pm$ in $F\setminus E$. The max and min functions are symmetric in their arguments, so conditions (b) and (c) are also well-defined on these pairs.

    Clearly, at most one of $w_1^\pm\ll w_2^\pm$ and $w_2^\pm \ll w_1^\pm$ can hold. Suppose neither hold. WLOG $\min\{L(w_1),L(w_1^{-1})\}=L(w_1)$, and $\min\{L(w_2),L(w_2^{-1})\}=L(w_2)$. It follows that $L(w_1)=L(w_2)$, and $L(w_1^{-1})=L(w_2^{-1})$. Thus the first and second halves of $w_1$ are equal to the first and second halves of $w_2$, hence $w_1=w_2$.

    Suppose that $w_1^\pm\ll w_2^\pm$ and $w_2^\pm\ll w_3^\pm$ hold. Then $l(w_1)\leq l(w_2)\leq l(w_3)$. If any one of these inequalities is strict, then $l(w_1)<l(w_3)$ and so $w_1\ll w_3$. Suppose then that $l(w_1)=l(w_2)=l(w_3)$. WLOG $\min\{L(w_i),L(w_i^{-1})\}=L(w_i)$.

    Then $L(w_1)\leq L(w_2)\leq L(w_3)$. If any one of these inequalities is strict, then $L(w_1)<L(w_3)$ and so $w_1\ll w_3$. Suppose then that $L(w_1)=L(w_2)=L(w_3)$. Then $L(w_1^{-1})<L(w_2^{-1})<L(w_3^{-1})$, so $L(w_1^{-1})<L(w_3^{-1})$, hence $w_1\ll w_3$.

    It follows that $\ll$ is a total order. Let $S$ be a non-empty subset of the pairs $w^\pm$ in $F\setminus E$. Let $S_a$ denote the subset of $S$ consisting of the pairs $w^\pm$ in $S$ with minimum length. Let $S_b$ denote the subset of $S_a$ consisting of the pairs $w^\pm$ in $S_a$, with $\min\{L(w),L(w^{-1})\}$ minimal. Finally let $S_c$ denote the subset of $S_b$ consisting of the pairs $w^\pm$ in $S_b$ with $\max\{L(w),L(w^{-1})\}$ minimal. Then $S_a,S_b,S_c\neq\emptyset$.

    It is clear that for all $w_1^\pm\in S_c$ and $w_2^\pm\in S\setminus S_c$ we have $w_1^\pm\ll w_2^\pm$. Also for all $w_1^\pm,w_2^\pm\in S_c$, neither $w_1^\pm\ll w_2^\pm$ nor $w_2^\pm\ll w_1^\pm$ holds, which implies $w_1^\pm=w_2^\pm$. Thus $S_c$ consists of a single element, which is the minimal element of $S$. Hence $\ll$ is a well-ordering.
  \end{solution}

\question Let $H$ be the subgroup $\overline{\{x^3, y^2, x^{-1}y^{-1}xy\}}$ of $F(\{x,y\})$ studied in Chapter 2. Find Schreier generators $B$ of $H$ using the Schreier transversal $\{e,x,x^2,y,xy,x^2y\}$. Describe a Nielsen transformation $\tau$ such that $B\tau$ is $N$-reduced.
  \begin{solution}
    We have that
    \begin{align*}
      B &= \{ uz\overline{uz}^{-1} \mid u\in U, z\in \{x,y\}, uz\notin U \} \\
        &= \{ x^3, yxy^{-1}x^{-1}, y^2, xyxy^{-1}x^{-2}, xy^2x^{-1}, x^2yxy^{-1}, x^2y^2x^{-2} \}.
    \end{align*}
    We think of $B$ having the ordering above. Note that (N2) fails for $B$, in particular $l((xyxy^{-1}x^{-2})x^3)=l(xyxy^{-1}x)<l(xyxy^{-1}x^{-2})$. Thus replacing $B$ by $B(41)$ fixes this issue. Still $l((x^2yxy^{-1})^{-1}x^3)=l(yx^{-1}y^{-1}x)<l(x^2yxy^{-1})$, so replacing $B(41)$ by $B(41)(6')(61)$ fixes this issue. Finally $l((x^2y^2x^{-2})x^3)=l(x^2y^2x)<l(x^2y^2x^{-2})$, so we should replace $B(41)(6')(61)$ with $B(41)(6')(61)(71)$. But still, $l((x^2y^2x)^{-1}x^3)=l(x^{-1}y^{-2}x)<l(x^2y^2x)$, so replacing what we have with $B(41)(6')(61)(71)(7')(71)$ will fix this.

    This leaves us with
    \begin{align*}
      B' &\coloneqq B(41)(6')(61)(71)(7')(71) \\
          &= (x^3, yxy^{-1}x^{-1}, y^2, xyxy^{-1}x^{-2}, xy^2x^{-1}, x^2yxy^{-1}, x^2y^2x^{-2})(41)(6')(61)(71)(7')(71) \\
          &= (x^3, yxy^{-1}x^{-1}, y^2, xyxy^{-1}x, xy^2x^{-1}, x^2yxy^{-1}, x^2y^2x^{-2})(6')(61)(71)(7')(71) \\
          &= (x^3, yxy^{-1}x^{-1}, y^2, xyxy^{-1}x, xy^2x^{-1}, yx^{-1}y^{-1}x^{-2}, x^2y^2x^{-2})(61)(71)(7')(71) \\
          &= (x^3, yxy^{-1}x^{-1}, y^2, xyxy^{-1}x, xy^2x^{-1}, yx^{-1}y^{-1}x, x^2y^2x^{-2})(71)(7')(71) \\
          &= (x^3, yxy^{-1}x^{-1}, y^2, xyxy^{-1}x, xy^2x^{-1}, yx^{-1}y^{-1}x, x^2y^2x)(7')(71) \\
          &= (x^3, yxy^{-1}x^{-1}, y^2, xyxy^{-1}x, xy^2x^{-1}, yx^{-1}y^{-1}x, x^{-1}y^{-2}x^{-2})(71) \\
          &= (x^3, yxy^{-1}x^{-1}, y^2, xyxy^{-1}x, xy^2x^{-1}, yx^{-1}y^{-1}x, x^{-1}y^{-2}x)
    \end{align*}
    We can see that $B'$ satisfies (N2), and also (N1). The only places in which (N3) can fail in $B_1$ are if we find $a,b,c\in B_1^\pm$ with $a\neq b^{-1}\neq c$, and $l(abc)\leq l(a)-l(b)+l(c)$, which can only occur if exactly half of $b$ and is cancelled in forming both $ab$ and $bc$.

    No such cancellation can occur. Thus (N3) also holds in $B'$, and thus $\tau=(41)(6')(61)(71)(7')(71)$ is the required Nielsen transformation.
  \end{solution}

  \stepcounter{question}

% \question Let $U$ be a Schreier transversal for a subgroup $H$ of $F(X)$ such that $u$ is of minimal length in $Hu$, for all $u\in U$. Prove that the Schreier basis $B$ is $N$-reduced.

\question Let $\phi\colon F(X)\twoheadrightarrow F(Y)$ be an epimorphism of free groups with $\abs{X}$ finite. Prove that $\abs{Y}$ is finite.
  \begin{solution}
    Since $F(X)=\langle X \rangle$, and $\phi$ is onto, it follows that $F(Y)=\langle X\phi \rangle$. Let $\tau$ be a Nielsen transformation such that $U=(X\phi)\tau$ is N-reduced. Then $U$ is a basis for $F(Y)$, and $\abs{U}\leq\abs{X\phi}\leq\abs{X}$ is finite. Since $Y$ is also basis for $F(Y)$, it follows that $\abs{Y}=\abs{U}$ is finite.
  \end{solution}

\question With $\phi$ as in the previous exercise, prove that $F(X)$ has a basis $V=V_1\cup V_2$ such that $\phi$ maps $V_1$ one-to-one onto a basis for $F(Y)$ and $V_2$ into $E$. (\emph{Hint:} Let $\tau$ be a Nielsen transformation carrying $X\phi$ into $Y$, and examine the effect of $\tau$ on $X$.)
  \begin{solution}
    Let $\tau$ be as in the previous exercise. We can think of $\tau$ as acting on $X$ in the following way. Suppose that $X=(x_1,\ldots,x_n)$ and $X\phi=(y_1,\ldots,y_m)$. When we encounter a (T1) transformation $(i')$ in $\tau$, we apply $(k')$ to the current $X$, for all $k$ such that $x_k\phi=y_i$. Note all such $(k')$ commute.

    Similarly when we encounter a (T2) transformation $(ij)$, we apply $(k\ell)$ to the current $X$, for all $k$ and for some $\ell$ such that $x_k\phi=y_i$ and $x_\ell\phi=y_j$. Note all such $(k\ell)$ commute. For a (T0) transformation we do nothing, but collect all affected entries into a set $V_2$ at the end, and all remaining members into a set $V_1$.

    If we set $V=X\tau$, where $\tau$ acting on $X$ is described above, then $V=V_1\cup V_2$. It is clear from the definitions of elementary Nielsen transformations, that $(V_1\tau)\phi=(X\phi)\tau$. Also $v\phi=e$ for all $v\in V_2$. Since $\tau$ acting on $X$ consists of only applying (T1) and (T2) Nielsen transformations, it follows that $V$ is a basis for $F(X)$.
  \end{solution}

\question Prove that if a free group $F$ contains an ascending chain of subgroups
  \[ H_1\subseteq H_2\subseteq\cdots\subseteq H_n\subseteq H_{n+1}\subseteq\cdots \]
  each characteristic in the next, then $H_i=H_j$ for all $i,j\geq$ some $N\in\N$.
  \begin{solution}
    Let $l_n=\min\{l(w) \mid e\neq w\in H_n\}$. Then $l_1\geq l_2\geq l_3\geq\cdots$. Since $l_n\in\N$, follows that for some $N\in\N$, $l_i=l_N$ for all $i\geq N$. Let $w\in H_N$ satisfy $l(w)=l_N$. We will show that $H_i=H_{i+1}$ for each $i\geq N$.

    Suppose $i\geq N$, and let $w\in H_i$ satisfy $l(w)=l_i=l_{i+1}$. Then there is a basis $B$ of $H_{i+1}$ containing $w$. For each $b\in B$, the automorphism of $H_{i+1}$ determined by interchanging $w$ and $b$ must fix $H_i$. This proves that $B\subset H_i$. Since $H_{i+1}=\langle B \rangle$, we have $H_i=H_{i+1}$. Thus $H_i=H_j$ for all $i,j\geq N$.
  \end{solution}

\question Let $\phi$ be an automorphism of $F(\{x,y\})$ and let $c=x^{-1}y^{-1}xy$. Prove that $c\phi=w^{-1}c^{\pm1}w$ for some $w\in F$.
  \begin{solution}
    For a regular Nielsen transformation $\tau$, let $\phi_\tau$ denote the homomorphism of $F(\{x,y\})$ which extends $\tau$. Then since $X\tau$ is a basis, $\phi_\tau$ is an automorphism.
    \begin{align*}
      c\phi_{(1)} &= (x^{-1}y^{-1}xy)\phi_{(1)}=xy^{-1}x^{-1}y=xc^{-1}x^{-1} \\
      c\phi_{(2)} &= (x^{-1}y^{-1}xy)\phi_{(2)}=x^{-1}yxy^{-1}=yc^{-1}y^{-1} \\
      c\phi_{(12)} &= (x^{-1}y^{-1}xy)\phi_{(12)}=(xy)^{-1}y^{-1}(xy)y=y^{-1}cy \\
      c\phi_{(21)} &= (x^{-1}y^{-1}xy)\phi_{(21)}=x^{-1}(yx)^{-1}x(yx)=x^{-1}cx.
    \end{align*}
    Therefore if $\tau$ is an elementary Nielsen transformation of type (T1) or (T2), then $c\phi_\tau=w^{-1}c^{\pm1}w$ for some $w\in F$. The restriction of $\phi$ to $X$ is some regular Nielsen transformation $\tau$. Write $\tau=\tau_1\cdots\tau_k$, where $\tau_i$ is an elementary Nielsen transformation of type (T1) or (T2) for each $1\leq i\leq k$.

    Since $\phi$ and $\phi_{\tau_1}\cdots\phi_{\tau_k}$ agree on $X$, they are equal as automorphisms. We know that $c\phi_{\tau_i}=w_i^{-1}c^{\pm1}w_i$ for some $w_i\in F$. Then for any $w\in F$, we have
    \[ (w^{-1}c^{\pm1}w)\phi_{\tau_i} = (w\phi_{\tau_i})^{-1}(w_i^{-1}c^{\pm1}w_i)^{\pm1}(w\phi_{\tau_i}) = (w_i(w\phi_{\tau_i}))^{-1}c^{\pm1}(w_i(w\phi_{\tau_i})). \]
    By induction, it follows that $c\phi=c\phi_{\tau_1}\cdots\phi_{\tau_k}=w^{-1}c^{\pm1}w$ for some $w\in F$.
  \end{solution}

\question Let $F$ be free of finite rank $r$. Then it was shown in deriving Corollary 7 that $\Aut F$ is generated by $r^2$ elements. Can you improve on this?
  \begin{solution}
    The elements of $\Aut F$ correspond to the regular Nielsen transformations of $X$. When $r=1$, $\abs{\Aut F}=2$, and so any set which generates $\Aut F$ has size at least 1. For $r=2$, let $\tau$ be the transposition given by $(u_1,u_2)\mapsto(u_2,u_1)$.

    By \textbf{1}, $\tau$ is a regular Nielsen transformation. Moreover, $(2')=\tau(1')\tau$ and $(21)=\tau(12)\tau$. Thus the regular Nielsen transformations are generated by $(1'),(12),\tau$, showing that $\Aut F$ can be generated by 3 elements in this case.

    When $r>2$, we will prove that $\Aut F$ is generated by 4 elements. Let $\tau$ be the transposition $(u_1,u_2,\ldots,u_r)\mapsto(u_2,u_1,\ldots,u_r)$, and let $\sigma$ be the $r$-cycle $(u_1,u_2,\ldots,u_r)\mapsto(u_2,\ldots,u_r,u_1)$. By \textbf{1}, $\tau$ and $\sigma$ are regular Nielsen transformations. Moreover they generate the group of permutations of $U$.

    For any $i,j$, let $\tau_{i,j}$ be a permutation of $U$ which puts $u_i$ and $u_j$ in the first and second positions respectively. Then $(ij)=\tau_{i,j}(12)\tau_{i,j}^{-1}$, and $(i')=\tau_{i,j}(1')\tau_{i,j}^{-1}$. Thus $(12)$, $(1')$, $\tau$ and $\sigma$ generate the regular Nielsen transformations, hence $\Aut F$ can be generated by 4 elements in this case.
  \end{solution}

\question Prove that $x^5y^{-1}x^8$, $xyxyx$ are primitive elements of $F(\{x,y\})$, but that $x^2$, $x^{-1}y^{-1}xy$ are not.
  \begin{solution}
    We have
    \begin{align*}
      (x,y)(1)(21)^5(1)(2)(21)^8 &= (x^{-1},y)(21)^5(1)(2)(21)^8 \\
                                 &= (x^{-1},yx^{-5})(1)(2)(21)^8 \\
                                 &= (x,x^5y^{-1})(21)^8 \\
                                 &= (x,x^5y^{-1}x^8),
    \end{align*}
    and
    \[ (x,y)(21)(12)(12) = (x,yx)(12)(12) = (xyxyx,yx). \]
    Thus $\{x,x^5y^{-1}x^8\}$ and $\{xyxyx,yx\}$ are both bases of $F(\{x,y\})$, and hence $x^5y^{-1}x^8$ and $xyxyx$ are primitive elements. If $\{x',y'\}$ is a basis of $F(\{x,y\})$, then the homomorphism determined by $x'\mapsto x$ and $y'\mapsto y$ must be an automorphism.

    Let $\phi$ be an automorphism of $F(\{x,y\})$. Then $x^2\phi=(x\phi)^2$. Since $x\phi\neq e$, we must have $l(x^2\phi)=l((x\phi)^2)\geq2$. Therefore $x^2\phi\neq x$.

    By \textbf{11} we know that if $c=x^{-1}y^{-1}xy$, then $c\phi$ is a conjugate of either $c$ or $c^{-1}$. Since $c$ is cyclically reduced, by \textbf{1.12}, this implies that $\hat{(c\phi)}$ is a cyclic rearrangement of $c$ or $c^{-1}$. Thus $l(c\phi)\geq l(c)=4$, and therefore $c\phi\neq x$.

     It follows that neither $x^2$ nor $x^{-1}y^{-1}xy$ belong to a basis of $F(\{x,y\})$, and hence they are not primitive elements.
  \end{solution}

\question Use Corollary 10 to prove that, for any set $X$, $F(X)$ is residually finite.
  \begin{solution}
    Let $e\neq w=x_1\cdots x_n\in F(X)$ be a reduced word in $X^\pm$. Let $Y=\{x_1^\pm,\ldots,x_n^\pm\}\cap X$. Then $Y$ is finite, and $w=x_1\cdots x_n$ is a reduced word in $Y^\pm$. By Corollary 10, there is a homomorphism $\phi\colon F(Y)\to G$, where $G$ is finite and $w\phi\neq1$. Define $\theta\colon X\to G$ by
    \[ x\theta =
      \begin{cases}
        x\phi & \text{if $x\in Y$,} \\
        1 & \text{otherwise.}
      \end{cases}
    \]
    This extends to a homomorphism $\theta'\colon F(X)\to G$, which agrees with $\phi$ on $F(Y)$. In particular $w\theta'\neq1$. Thus $F(X)$ is residually finite.
  \end{solution}

\question Let $H$ be a subgroup of finite index in a group $G$. Prove that $H$ is residually finite if and only if $G$ is.
  \begin{solution}
    Suppose $G$ is residually finite, and let $e\neq h\in H$. Then there is a homomorphism $\phi\colon G\to F$, with $F$ finite and $h\phi\neq1$. Then $\phi{\restriction_H}\colon H\to F$ satisfies $h\phi{\restriction_H}=h\phi\neq1$. Therefore $H$ is residually finite.

    Conversely, suppose that $H$ is residually finite, and let $e\neq g\in G$. If $g\notin H$, then since $G$ acts on the set of cosets $G/H$, this induces a homomorphism $\phi\colon G\to S_n$, where $n=\abs{G:H}$. Clearly $\ker\phi\leq H$, so $g\phi\neq1$, and $G/\ker\theta$ is finite.

    If $g\in H$, then there exists a homomorphism $\phi\colon H\to F$ with $g\phi\neq1$, for some finite group $F$. Now $\abs{H:\ker\phi}$ and $\abs{G:H}$ are finite, thus $\abs{G:\ker\phi}$ is finite. Since $G$ acts on the set of cosets $G/\ker\phi$, this induces a homomorphism $\phi'\colon G\to S_m$, where $m=\abs{G:\ker\phi}$. Again, $\ker\phi'\leq\ker\theta$, so $g\phi'\neq1$, and $G/\ker\phi'$ is finite.

    It follows that $G$ is residually finite, completing the proof.
  \end{solution}
\end{questions}

%%% Local Variables:
%%% mode: latex
%%% TeX-master: "johnson"
%%% End:
\section{Free Presentations of Groups}
\begin{questions}
\question Give an alternative proof of Proposition 1 by showing that if $G$ is any group, then $G$ has the presentation $\langle X \mid M \rangle$, where $X=G$ and
  \[ M = \{xy = (x,y)m \mid x,y\in G\} \]
  is the multiplication table of $G$ (where $(x,y)m$ stands in the $(x,y)$-place).

\question Let $G=\langle X \mid R \rangle$ and $G'$ be the commutator subgroup of $G$. Prove that $G_{ab}=G/G'$ has the presentation
  \[ \langle X \mid R, [X,X] \rangle, \]
  where $[X,X]=\langle [x,y] \mid x,y\in X,\ x\neq y \rangle$. Use Proposition 3 to confirm the result of Exercise 1.7.

\question For $l,m,n\in\Z$, prove that
  \[ D(l,m,n) \cong D(m,l,n) \cong D(-l,m,n), \]
  and deduce that $D(l,m,n)$ is independent of the orders and signs of $l,m,n$.

\question Identify $D(l,m,n)$ in the following cases:
  \begin{parts}
  \part one of $l,m.n$ is 1,
  \part two of $l,m,n$ are 2.
  \end{parts}

\question Let $D_\infty=\langle a,b \mid a^2, b^2 \rangle$. Identify $D_\infty/D_\infty'$. Prove that $D_\infty'$ is a) cyclic, and b) infinite. (Hint: For a), consider conjugates of $[a,b]$ and for b), consider homomorphisms from $D_\infty$ onto the groups $D(2,2,n)$.)

\question Show that the groups
  \[ \langle a,b \mid a^bb^a = (b^{-1}a^2)^2 = e \rangle, \quad \langle x,y \mid x^2 = y^3 = e \rangle \]
  are isomorphic.

\question Do you recognize the group
  \[ \langle a,b,c \mid a^3=b^2=c^2=(ab)^2=(bc)^2=[a,c]=e \rangle\ ? \]
  What is the order of the element $abc$?

\question Prove the converse of the result in (10) of Example 7:
  \[ (m,n)\neq1 \implies Z_{mn} \not\cong Z_m\times Z_n. \]

\question Prove that the groups
  \[ \langle x,y \mid xyx = y, yxy = x \rangle, \quad \langle a,b \mid a^2=b^2, a^{-1}ba=b^{-1} \rangle \]
  are isomorphic.

\question Show that the groups
  \[ \langle a,b,c,d \mid ab=c, bc=d, cd=a, da=b \rangle \]
  \[ \langle a,b,c,d,f,g \mid ab=d, bc=f, cd=g, df=a, fg=b, ga=c \rangle \]
  are cyclic and find their orders.

\question Consider the group $F(2,6)$ given by
  \[ \langle a,b,c,d,f,g \mid ab=c, bc=d, cd=f, df=g, fg=a, ga=b \rangle. \]
  Prove that there is a homomorphism $\chi\colon F(2,6)\to\Sym(\Z)$ such that
  \[ n(a\chi)=n+1,\ n(b\chi) = -n,\ n\in\Z, \]
  and deduce that $F(2,6)$ is an infinite group.

\question Prove that the group $F(2,5)$ given by
  \[ \langle a,b,c,d,f \mid ab=c, bc=d, cd=f, df=a, fa=b \rangle \]
  is a finite cyclic group.

\question Check the \emph{Witt Identity}
  \[ [[x,y],z^x][[z,x],y^z][[y,z],x^y] = e, \]
  and use it to prove that the group
  \[ \langle x,y,z \mid [x,y]=z, [y,z]=x, [z,x]=y \rangle \]
  is trivial.

\question Let $x$ and $y$ be members of a group such that $x^{-1}y^nx=y^m$, for some $m,n\in\Z$. Prove that, for any $k\in\N$, $x^{-1}y^{n^k}x=y^{m^k}$. Deduce that the group
  \[ \langle x,y \mid y^{-1}x^ny = x^{n+1}, x^{-1}y^nx = y^{n+1} \rangle \]
  is trivial for any $n\in\Z$.

\question Given $G=\langle X\mid R \rangle$, prove that
  \begin{parts}
  \part if $G$ is finitely generated, then there is a finite subset $X'\subseteq X$ such that $G=\langle X' \rangle$,
  \part if $G$ is in addition finitely presented, then there is in addition, a finite subset $R'$ of $F(X')$ such that $G=\langle X' \mid R' \rangle$.
  \end{parts}

\question Solve the word problem for finite groups according to the following scheme:
  \begin{parts}
  \part let $G=\langle X\mid R \rangle$ be a finite presentation of a finite group $G$, and let $w\in F(X)$;
  \part enumerate the elements of $\overline{R}\lhd F(X)$ as products of conjugates of $R^\pm$;
  \part enumerate the homomorphisms: $G\to S_n$, $n\in\N$ as mappings $X\to S_n$ that vanish on $R$, and
  \part by examining alternate entries in the lists (b) and (c), decide in a finite number of steps whether $w=e$ in $G$ or not.
  \end{parts}

\question Let $G$ be the group
  \[ \langle x,y \mid x^2y=y^2x, x^8=e \rangle. \]
  Prove that $y^8=e$ in $G$ and draw the associated van Kampen diagram.

\question Draw a van Kampen diagram for the group
  \[ \langle x,y \mid x^3=y^2, x^3=e, y^{-1}xy=x^2 \rangle, \]
  which shows that the second relation is superfluous.

\question Draw a van Kampen digram which shows that the group
  \[ \langle a,b \mid abab^2 = e = baba^2 \rangle \]
  is abelian.
\end{questions}

%%% Local Variables:
%%% mode: latex
%%% TeX-master: "johnson"
%%% End:

\outersection{Computation with Finitely Generated Abelian Groups}

%% Sections numbered with X
%% Reset section numbering
\renewcommand{\thesection}{\arabic{section}}
\setcounter{section}{0}

\begin{abstract}
  The aim of this part is to better understand the theory of finitely generated abelian groups, and computational methods pertaining to them. We shall see that an important tool is the \emph{Smith Normal Form} of an integer matrix.

Naive methods inspired by Gaussian elimination to compute the Smith Normal Form of a matrix can require arbitrary precision arithmetic for small inputs. We will explore some useful techniques, which allow us to perform calculations with respect to an appropriate modulus.
\end{abstract}

\section{Preliminaries}

In these notes, we will always write the group operation of an abelian group additively, unless otherwise stated. 

\begin{definition}[Universal Property]
  A group $G$ is \emph{free abelian} on a subset $X\subseteq G$ if every map from $X$ to an abelian group $H$ extends to a unique homomorphism $G\to H$. We call $X$ a \emph{basis} for $G$ if $G$ is free abelian on $X$.
\end{definition}

There is a similarity with vector spaces: $B$ is a basis of a vector space $V$ if and only if every map from $B$ to a vector space $W$ extends uniquely to a linear map $V\to W$. The existence of such a linear map is equivalent to $B$ being a linearly independent set, and the uniqueness is equivalent to $B$ spanning $V$.


\begin{theorem}
  The free abelian groups with finite basis, up to isomorphism, consist of the groups $\Z^n$ for $n\in\N$.
\end{theorem}

\begin{proof}
  Let $e_i\in\Z^n$ be the $i$th standard basis vector. Each element of $\Z^n$ takes the form $m_1e_1+\cdots m_ne_n$, for unique integers $m_i\in\Z$. Given a map from $\{e_1,\ldots,e_n\}$ to a group $H$, we can extend it to a group homomorphism $\Z^n\to H$ as follows. If $e_i\mapsto h_i$, then define
  \[ m_1e_1 + \cdots + m_ne_n \mapsto m_1h_1 + \cdots + m_nh_n. \]
  It is readily seen that this defines a group homomorphism from $\Z^n$ to $H$, sending $e_i\mapsto h_i$. Moreover each group homomorphism $\Z^n\to H$ which maps $e_i\mapsto h_i$ must agree with this. Hence $\Z^n$ is free abelian on $\{e_1,\ldots,e_n\}$.

  Conversely, suppose that $G$ is free abelian on $X=\{x_1,\ldots,x_n\}$. By the universal property, we obtain a homomorphism $\phi\colon G\to\Z^n$ which maps $x_i\mapsto e_i$. Similarly we have a homomorphism $\psi\colon\Z^n\to G$ which maps $e_i\mapsto x_i$. Then $\psi\circ\phi$ is an endomorphism of $G$ fixing $X$. By the uniqueness in the universal property, $\psi\circ\phi$ must be the identity on $G$. Similarly $\phi\circ\psi$ is the identity on $\Z^n$. Hence $G\cong\Z^n$.
\end{proof}

Let $G$ be an abelian group generated by $n$ elements. Since $\Z^n$ is free abelian there is an epimorphism $\Z^n\twoheadrightarrow G$ with kernel $H$. By the first isomorphism theorem, $\Z^n/H$ is isomorphic to $G$. Thus understanding the structure of subgroups of $\Z^n$ will give us insight into the structure of finitely generated abelian groups.

\begin{theorem}[Dedekind]
  A subgroup of $\Z^n$ can be generated by at most $n$ elements.
\end{theorem}

\begin{proof}
  We proceed by induction on $n$. Defining $\Z^0\coloneqq\{0\}$, the case $n=0$ becomes trivial. Let $n>0$, and let $\varphi\colon\Z^n\twoheadrightarrow\Z^n/\langle e_n \rangle$ be the natural map. Note that $\Z^n/\langle e_n \rangle\cong\Z^{n-1}$. If $H\leq\Z^n$, then $\varphi(H)$ is isomorphic to a subgroup of $\Z^{n-1}$. Inductively, we may assume $\varphi(H)=\langle h_1 + \langle e_n \rangle, \ldots, h_{n-1} + \langle e_n \rangle \rangle$, for $h_i\in H$. Note that $H\cap\langle e_n \rangle$ is cyclic, so let $H\cap\langle e_n \rangle=\langle h_n \rangle$ for some $h_n\in H$. We claim $H=\langle h_1,\ldots,h_n \rangle$.

  If $h\in H$, then $\varphi(h)=h'+\langle e_n \rangle$ for some $h'\in\langle h_1,\ldots,h_{n-1} \rangle$. Now $h-h'\in H\cap\langle e_n \rangle=\langle h_n \rangle$. Thus $h\in\langle h_1,\ldots,h_n \rangle$, and hence $H=\langle h_1,\ldots,h_n \rangle$. By induction the theorem follows.
\end{proof}

\begin{definition}
  The \emph{integer row space} of an $m\times n$ integer matrix $A$ is $S(A)\coloneqq\{xA : x\in\Z^m\}$, i.e. the set of all integral linear combinations of rows of $A$.
\end{definition}

If $H=\langle h_1,\ldots,h_m \rangle\leq\Z^n$, then $H$ consists of all integral linear combinations of the $h_i$. Hence each finitely generated abelian group is isomorphic to $\Z^n/S(A)$ for some $m\times n$ integer matrix $A$. For most descriptions of finitely generated abelian groups, we can explicitly find such an $A$. For example, consider the abelianisation $G_{ab}\coloneqq G/[G,G]$ of a finitely presented group $G$.

If $G=\langle X \mid R \rangle$ is a finite presentation, then the abelianisation $G_{ab}$ is a finitely generated abelian group, with presentation $\langle X \mid R, [X,X] \rangle$. If $X=\{x_1,\ldots,x_n\}$, then we can rewrite each of the relations in $R$ in the form $x_1^{\alpha_1}\cdots x_n^{\alpha_n}$, $\alpha_i\in\Z$. The epimorphism $\Z^n\twoheadrightarrow G_{ab}$ sending $e_i\mapsto x_i$ has kernel generated by the $(\alpha_1,\ldots,\alpha_n)$ appearing in the relations. Using these rows we can form a ``relation matrix'' $A$, and then $G_{ab}\cong\Z^n/S(A)$.

\section{Integer Row Reduction}

A notion of row space can be defined for matrices over an arbitrary ring. When working over a field, we can apply row operations to produce a matrix in reduced row echelon form (e.g. using Gaussian elimination). We develop a similar theory for integer matrices below.

\begin{definition}
  An \emph{integer row operation} applied to a matrix is one of the following:
  \begin{enumerate}
  \item Swap two rows.
  \item Multiply a row by $-1$.
  \item Add an integer multiple of one row to another row.
  \end{enumerate}
  Two $m\times n$ integer matrices $A$ and $B$ are \emph{row equivalent} if there is a sequence of integer row operations transforming one into the other, and we write $A\sim B$.
\end{definition}

We can reverse each integer row operation with another. The first two types are involutions. To reverse adding $q$ times row $i$ to row $j$, for $i\neq j$, we add $-q$ times row $i$ to row $j$. This makes $\sim$ an equivalence relation on integer $m\times n$ matrices.

In the above definition rows may only be multiplied by $-1$, in contrast to row operations over a field, where rows may be multiplied by an arbitrary non-zero scalar. This is because multiplying a row by a scalar can be reversed when multiplying by a unit, and the units of a field are its non-zero elements, but the units of $\Z$ are just $1$ and $-1$.

\begin{theorem}
  If $A$ and $B$ are integer matrices with $A\sim B$, then $S(A)=S(B)$.
\end{theorem}

\begin{proof}
  If $B$ is obtained from $A$ by a single integer row operation, then the rows of $B$ are clearly in $S(A)$, so $S(B)\subseteq S(A)$. Moreover this row operation can be reversed, thus $S(A)\subseteq S(B)$. Therefore $A\sim B\implies S(A)=S(B)$.
\end{proof}

The corresponding notion of reduced row echelon form for integer matrices is row Hermite Normal Form (HNF).

\begin{definition}
  An integer $m\times n$ matrix $A$ is in \emph{row Hermite Normal Form (HNF)} if
  \begin{enumerate}
  \item The nonzero rows of $A$ are the first $r$ rows of $A$, for some $r\leq m$.
  \item If $j_i$ is minimal with $A_{i,j_i}$ nonzero, for $1\leq i\leq r$, then $j_1<j_2<\cdots<j_r$.
  \item If $1\leq i\leq r$, then $A_{i,j_i}>0$.
  \item If $1\leq k<i\leq r$, then $0\leq A_{k,j_i}<A_{i,j_i}$.
  \end{enumerate}
\end{definition}

In the above definition, the entries $A_{i,j_i}$ behave similarly to the pivot entries in reduced row echelon form. Below is a sample matrix in row Hermite Normal Form.
\[ A \coloneqq
  \begin{bmatrix}
    2 & 1 & 2 & 3 \\
    0 & 0 & 7 & 5 \\
    0 & 0 & 0 & 9
  \end{bmatrix}.
\]

Suppose we want to determine whether a given $u\in\Z^4$ is in $S(A)$. For instance, let $u=(4,2,-3,10)$. We seek $x,y,z\in\Z$ with $u=xa_1+ya_2+za_3$, where $a_1,a_2,a_3$ are the rows of $A$. Isolating the first column, we see $4=2x$, so $x=2$. Next set $v=u-2a_1=(0,0,-7,4)$. We require $v=ya_2+za_3$. The second column holds no information. The third column tells us that $-7=7y$, so $y=-1$. Set $w=v+a_2=(0,0,0,9)$. We require $w=za_3$. Observing the last column shows $9=9z$, so $z=1$. Hence $u=2a_1-a_2+a_3$, and so $u\in S(A)$.

If at any stage we found an equation that was not solvable for integer $x,y,z$, then we would instead conclude that $u\notin S(A)$. Clearly, testing membership in $S(A)$ is easy when $A$ is in HNF. One can also show the nonzero rows of $A$ form a basis of $S(A)$ when $A$ is in HNF.

\begin{theorem}
  If $A$ is an $m\times n$ integer matrix, then there is a unique $m\times n$ integer matrix $B$ with $A\sim B$ and $B$ in row Hermite Normal Form.
\end{theorem}

\begin{proof}
  We prove this by induction. The result holds trivially when $m=0$ or $n=0$. Suppose that $m,n\geq1$, and that the result holds for all smaller matrices. If there are two nonzero entries in the first column, say $0<\abs{A_{k,1}}\leq\abs{A_{\ell.1}}$ with $k\neq\ell$, then we decrease the quantity $\abs{A_{1,1}}+\cdots+\abs{A_{m,1}}$ as follows.

  First multiply rows $k,\ell$ by $-1$ if necessary so that $A_{k,1},A_{\ell,1}>0$. Next, subtract row $k$ away from row $\ell$. Since $0\leq A_{\ell,1}-A_{k,1}<A_{\ell,1}$, this strictly decreases the quantity $\abs{A_{1,1}}+\cdots+\abs{A_{m,1}}$. This quantity is a non-negative integer, so it can only be decreased finitely many times. Hence we may assume that $A$ has at most one nonzero entry in the first column. If all entries in the first column are zero, then $A$ has the block form

  \[ A = \left[
      \begin{array}{c|ccccc}
        0 &&&&& \\
        \vdots &&& A' && \\
        0 &&&&& \\
      \end{array}\right]
  \]

  By induction we can reduce $A'$ to HNF by row operations, thus we can reduce $A$ to HNF by row operations. Suppose that $A$ has only one nonzero entry in the first column. By swapping rows and multiplying by $-1$ if necessary, we may assume that the nonzero entry is $A_{1,1}$ and $A_{1,1}>0$. Now $A$ has the block form
  \[ A = \left[
      \begin{array}{c|ccccc}
        A_{1,1} & A_{1,2} && \cdots && A_{1,n} \\
        \hline
        0 &&&&& \\
        \vdots &&& A' && \\
        0 &&&&& \\
      \end{array}\right]
  \]
  By induction we can reduce $A'$ to HNF by row operations, so we may assume that $A'$ is in HNF. Suppose that the nonzero rows of $A$ are now the first $r$ rows, and that the first nonzero entry in each row is given by $A_{i,j_i}$, for $1\leq i\leq r$. Let $2\leq k\leq r$, and suppose we have already arranged for $0\leq A_{1,j_i}<A_{i,j_i}$ to hold, for each $i=2,\ldots,k-1$.

  Using the division algorithm we may write $A_{1,j_k}=qA_{k,j_k}+r$, where $0\leq r<A_{k,j_k}$. We subtract $q$ times row $k$ away from row 1. Because $A_{k,j_k}$ is the first nonzero entry in row $k$, and $1=j_1<j_2<\cdots<j_k$, we still have $0\leq A_{1,j_i}<A_{i,j_i}$ for each $i=2,\ldots,k-1$.

  Hence by induction we can reduce $A$ to some matrix $B$ in HNF, by integer row operations. Let $B'$ be another matrix for which $A\sim B'$ and $B'$ is in HNF. Let $b_1,\ldots,b_m$ and $b_1',\ldots,b_m'$ be the rows of $B$ and $B'$ respectively. If $B\neq B'$, then there are entries with $B_{i,j}\neq B'_{i,j}$. Choose such $i,j$ with $j$ minimal, without loss of generality $B_{i,j}>B_{i,j}'$. We have $b_i,b_i'\in S(B)=S(A)=S(B')$, hence $b_i-b_i'\in S(B)$.

  Suppose that only the first $r$ rows of $B$ are nonzero, and let $B_{i,j_i}$ be the first nonzero entry in $b_i$ for $1\leq i\leq r$. The first $j-1$ entries of $b_i-b_i'$ are zero, so $b_i-b_i'$ is an integral linear combination of the rows $b_k$ with $j_k\geq j$. However $b_{ij}-b_{ij}'\neq0$, and so we must have $j_k=j$ for some $k$ and $B_{k,j}\mid B_{i,j}-B_{i,j}'$. Since $0\leq B_{i,j}'< B_{i,j}<B_{k,j}$, we must have $\abs{B_{i,j}-B_{i,j}'}<B_{k,j}$, so $B_{i,j}-B_{i,j}'=0$, which is a contradiction. Therefore $B=B'$, and so each integer matrix $A$ is row equivalent to a unique integer matrix $B$ in HNF.
\end{proof}

\begin{corollary}
  Let $A,B$ be integer matrices. Then $S(A)=S(B)\implies A\sim B$.
\end{corollary}

\begin{proof}
  It suffices to prove that if $A$ and $B$ are in HNF, then $S(A)=S(B)\implies A=B$. We refer the reader to \cite{sims} for the proof of this.
\end{proof}

The proof of this theorem is readily turned into an procedure, e.g. this one (from \cite{sims})

\begin{algorithm}[H]
  \SetKwInOut{Procedure}{Procedure}
  \Procedure{ROW\_REDUCE(\textasciitilde$A$);}
  \KwIn{An $m\times n$ integer matrix $A$}
  \KwResult{Integer row operations are applied to $A$ to reach row Hermite Normal Form}

  $A\coloneqq B$; $i\coloneqq1$; $j\coloneqq1$;

  \While{$i\leq m$ and $j\leq n$}{
    \eIf{$A_{k,j}=0$ for $i\leq k\leq m$}{$j\coloneqq j+1$}{
      \While{there exist distinct $k,\ell$ with $i\leq k,\ell\leq m$ and $0<\abs{A_{k,j}}\leq\abs{A_{\ell,j}}$}{
        $q\coloneqq A_{\ell,j} \bdiv A_{k,j}$\;
        Subtract $q$ times row $k$ of $A$ from row $\ell$
      }
      Let $A_{kj}\neq0$ with $i\leq k\leq m$; (this $k$ is unique)\\
      \If{$k\neq i$}{swap rows $i$ and $k$ of $A$}
      \If{$A_{i,j}<0$}{multiply row $i$ of $A$ by $-1$}
      \For{$\ell\coloneqq 1$ to $i-1$}{
        $q\coloneqq A_{\ell,j} \bdiv A_{i,j}$\;
        Subtract $q$ times row $i$ of $A$ from row $\ell$
      }
      $i\coloneqq i+1$; $j\coloneqq j+1$\;
    }
  }
\end{algorithm}

There is some freedom when implementing the above algorithm. In the inner while loop, we must select which indices $k,\ell\geq i$ are selected with $0<\abs{A_{k,j}}\leq\abs{A_{\ell,j}}$. If $k,\ell$ are chosen so that $\abs{A_{\ell,j}}-\abs{A_{k,j}}$ is maximized, then that the quantity $\abs{A_{1,j}}+\cdots+\abs{A_{m,j}}$ decreases as much as possible in each iteration. As described in the second paragraph of the previous proof, the condition in this while loop no longer holds when this quantity can no longer decrease. Thus one might expect this strategy to be the most efficient.

In practice one finds that this strategy can result in $A$ having large entries, significantly increasing run time. This issue is discussed in \cite{rosser}, where an alternate strategy is proposed. The \emph{Rosser strategy} is to choose $k,\ell$ so that $\abs{A_{\ell,j}}$ is as large as possible, and $\abs{A_{k,j}}$ is as large as possible with $k\neq\ell$. Many authors recommend this because it tends to keep the size of the entries small during the procedure. See \autoref{hnfcode} for a {\sc Magma} implementation of the row reduction algorithm employing this strategy.

\section{Smith Normal Form}
\label{snf-discussion}

In the previous section we showed that integer row operations applied to an $m\times n$ matrix $A$ do not change $S(A)$. We prove further that column operations do not affect the isomorphism type of $\Z^n/S(A)$.

\begin{definition}
  An \emph{integer column operation} applied to a matrix is one of the following:
  \begin{enumerate}
  \item Swap two columns.
  \item Multiply a row by $-1$.
  \item Add an integer multiple of one column to another column.
  \end{enumerate}
\end{definition}

\begin{definition}
  Two $m\times n$ integer matrices $A$ and $B$ are \emph{equivalent} if there is a sequence of integer row and column operations transforming one into the other, and we write $A\approx B$.
\end{definition}

\begin{theorem}
  If $A,B$ are $m\times n$ integer matrices, then $A\approx B$ if and only if there exist matrices $P\in\GL_m(\Z)$ and $Q\in\GL_n(\Z)$ such that $B=PAQ$.
\end{theorem}

\begin{proof}
  For each integer row operation on an $m\times n$ matrix, there is a corresponding matrix $P\in\GL_m(\Z)$ such that the effect of applying the row operation is equivalent to left multiplication by $P$. Moreover these ``elementary'' matrices generate $\GL_m(\Z)$. Similarly integer column operations correspond to right multiplication by matrices in $\GL_n(\Z)$.
\end{proof}

\begin{theorem}
  \label{commutative-diagram}
  Let $A,B$ be integer $m\times n$ matrices. Then $A\approx B\implies \Z^n/S(A)\cong\Z^n/S(B)$.
\end{theorem}

\begin{proof}
  There exist $P\in\GL_m(\Z)$, $Q\in\GL_n(\Z)$ such that $B=PAQ$. Then $B\sim AQ$, hence $S(B)=S(AQ)$. Notice that $S(A)Q\coloneqq\{uQ : u\in S(A)\}=\{xAQ : x\in\Z^m\}=S(AQ)$. It follows that the mapping illustrated in the diagram is a well defined homomorphism $\Z^n/S(A)\to\Z^n/S(AQ)=\Z^n/S(B)$. It is an isomorphism because $Q$ is invertible.

  \centering
  \begin{tikzcd}
    \Z^n \arrow[rrr, "Q\colon x\;\mapsto\; xQ"] \arrow[dd, two heads] & & & \Z^n \arrow[dd, two heads] \\
    & & &  \\
    \Z^n/S(A) \arrow[rrr, "x + S(A)\; \mapsto\;  xQ + S(AQ)"'] & & & \Z^n/S(AQ)
  \end{tikzcd}

\end{proof}

We establish a normal form which distinguishes the isomorphism types of $\Z^n/S(A)$.

\begin{definition}
  An $m\times n$ integer matrix $A$ is in \emph{Smith Normal Form (SNF)} if there is some $r$ such that $d_i\coloneqq A_{i,i}>0$ for each $1\leq i\leq r$, all remaining entries of $A$ are zero, and $d_1\mid d_2\mid \cdots \mid d_r$. The integers $d_1,\ldots,d_r$ are the \emph{invariant factors} of $A$.
\end{definition}

Here is a sample matrix in Smith Normal Form
\[ A \coloneqq
  \begin{bmatrix}
    2 & 0 & 0 & 0 & 0 \\
    0 & 4 & 0 & 0 & 0 \\
    0 & 0 & 12 & 0 & 0
  \end{bmatrix}.
\]
Note that $S(A)=\{(2x,4y,12z,0,0) : x,y,z\in\Z\}$. This is the kernel of the epimorphism
\begin{align*}
  \Z^5 &\twoheadrightarrow \Z_2\oplus\Z_4\oplus\Z_{12}\oplus\Z\oplus\Z \\
  (a,b,c,d,e) &\mapsto (a\bmod 2, b\bmod 4, c\bmod 12, d, e).
\end{align*}
Therefore $\Z^5/S(A)\cong\Z_2\oplus\Z_4\oplus\Z_{12}\oplus\Z\oplus\Z$. Determining the isomorphism type of $\Z^n/S(A)$ for an $m\times n$ matrix $A$ is simple when $A$ is in Smith Normal Form.

\begin{theorem}
  \label{decomposition}
  Let $A$ be an $m\times n$ integer matrix in SNF, with invariant factors $d_1,d_2,\ldots,d_r$. Then $\Z^n/S(A)\cong\Z_{d_1}\oplus\cdots\oplus\Z_{d_r}\oplus\Z^{n-r}$.
\end{theorem}

\begin{theorem}
  \label{SNF-existence}
  If $A$ is an integer matrix, then there exists an integer matrix $B$ with $A\approx B$ and $B$ in Smith Normal Form.
\end{theorem}

\begin{proof}
  Suppose $A$ is not the zero matrix. The first goal is to reduce $A$ to the block form
  \[ \left[
    \begin{array}{c|ccc}
      d & 0 & \cdots & 0 \\
      \hline
      0 & & & \\
      \vdots & & A' & \\
      0 & & & \\
    \end{array}\right] \tag{$\star$}
  \]
  with $d>0$. Let $A_{i,j}$ be a nonzero entry of $A$. Swap rows $i$ and 1, and columns $j$ and 1 if required, to ensure $i=j=1$, and multiply row 1 by $-1$ if necessary to ensure $A_{1,1}>0$. If $A_{1,1}$ divides all entries in row and column 1, then add multiples of row/column 1 to the other rows/columns, to reach the form ($\star$).

  If $A_{1,1}$ does not divide all entries in row and column 1, then we decrease this entry as follows. Suppose that $A_{1,1}\nmid A_{i,1}$ for some $i$. Write $A_{i,1}=qA_{1,1}+r$ with $0<r<A_{1,1}$. Then add $-q$ times row 1 to row $i$, and swap rows 1 and $i$. Similarly, if $A_{1,1}\nmid A_{1,j}$ for some $j$, then write $A_{1,j}=qA_{1,1}+r$ with $0<r<A_{1,1}$. Then add $-q$ times column 1 to column $j$, and swap columns 1 and $j$.

  Therefore, after repeating this enough times, we ensure $A_{1,1}$ divides all entries in row and column 1, and thus we can reach the block form ($\star$). Inductively, we can reduce $A'$ to SNF by row and column operations, so we can reduce $A$ to a matrix with block form
  \[ \left[
      \begin{array}{ccc|cc}
        d_1 &&&& \\
            & \ddots &&& 0 \\
            && d_r && \\
        \hline
            &&&& \\
        & 0 &&& 0 \\
      \end{array}\right]
  \]
  where $d_i>0$. If the divisibility condition $d_1\mid d_2\mid\cdots\mid d_r$ holds, then we have produced a matrix in SNF. If not, suppose that $d_i\nmid d_j$ for some $i<j$. Set $a\coloneqq d_i$ and $b\coloneqq d_j$. The Euclidean algorithm produces integers $u,v$ with $d\coloneqq\gcd(a,b)=ua+vb$, and then $\ell\coloneqq\lcm(a,b)=ab/d$. Now
  \[
    \begin{bmatrix}
      u & v \\
      -b/d & a/d
    \end{bmatrix}
    \begin{bmatrix}
      a & 0 \\
      0 & b
    \end{bmatrix}
    \begin{bmatrix}
      1 & -vb/d \\
      1 & ua/d
    \end{bmatrix}
    =
    \begin{bmatrix}
      d & 0 \\
      0 & \ell
    \end{bmatrix}.
  \]
  In this way, we can use row and column operations to replace $d_i$ and $d_j$ with $\gcd(d_i,d_j)$ and $\lcm(d_i,d_j)$ respectively. Repeating this will produce a matrix in Smith Normal Form.
\end{proof}

Uniqueness follows from the uniqueness of the integers in the following important theorem. Note that the previous result implies the existence of such integers. To prove the uniqueness, we refer the reader to \cite{dummit}.

\begin{theorem}[Fundamental Theorem of Finitely Generated Abelian Groups]
  \label{FToFGAG}
  Let $G$ be a finitely generated abelian group. Then there exists unique integers $n,d_1,\ldots,d_r>0$ with $d_1\mid d_2\mid\cdots\mid d_r$, and
  \[ G\cong\Z_{d_1}\oplus\cdots\oplus\Z_{d_r}\oplus\Z^{n-r}. \]
\end{theorem}

\begin{corollary}
  \label{unique-snf}
  If $A$ is an $m\times n$ integer matrix, then there is a unique $m\times n$ integer matrix $B$ with $A\approx B$ and $B$ in Smith Normal Form.
\end{corollary}

Thus it makes sense to define the \emph{invariant factors} of an integer matrix to be the invariant factors of the matrix equivalent to it which is in SNF.

\begin{corollary}
  If $A$ and $B$ are $m\times n$ integer matrices with $\Z^n/S(A)\cong\Z^n/S(B)$, then $A\approx B$.
\end{corollary}

\begin{proof}
  By \autoref{FToFGAG}, $A$ and $B$ have the same invariant factors. Because they have the same dimensions, they must be equivalent to the same matrix in SNF. Hence $A\approx B$.
\end{proof}

The proof of \autoref{SNF-existence} is readily transformed into a procedure for computing the Smith Normal Form of a given matrix. See \autoref{snfcode} for a {\sc Magma} implementation. However one quickly finds that it has certain defects. For example, I ran the implementation in \autoref{snfcode} on a random $100\times 100$ matrix with entries in $\{-1,0,1\}$, on the machine described in \autoref{conclusion}. After 13 hours the procedure had still not terminated. We investigate what happens with a smaller example. Consider the following matrix

\[
  A\coloneqq\begin{bsmallmatrix}
    6 & 9 & 8 & 2 & 6 & 5 & 2 & 0 \\
    5 &10 & 1 & 4 & 8 & 5 & 4 & 8 \\
    3 & 1 & 9 & 6 & 3 &10 & 5 & 3 \\
    6 &10 & 3 & 2 & 5 & 0 & 3 &10 \\
    6 & 1 & 3 & 8 & 6 & 6 & 9 & 1 \\
    6 & 2 & 2 & 1 & 8 & 1 &10 &10 \\
    4 & 4 &10 & 7 &10 & 8 & 3 & 3 \\
    7 & 0 & 1 & 9 & 8 & 5 & 5 & 1 \\
  \end{bsmallmatrix}
\]

If we run the {\sc Magma} implementation given in \autoref{snfcode} on $A$, we find that the intermediate matrix in which the entry with largest modulus appears, is
\[
  \begin{bsmallmatrix}
    1 & 0 & 0 & 0 & 0 & 0 & 0 & 0 \\
    0 & 1 & 0 & 0 & 0 & 0 & 0 & 0 \\
    0 & 0 & 1 & 0 & 0 & 0 & 0 & 0 \\
    0 & 0 & 0 & 1 & 0 & 0 & 0 & 0 \\
    0 & 0 & 0 & 0 & 1 & 0 & 0 & 0 \\
    0 & 0 & 0 & 0 & 0 & 1 & 48829330326663043031960 & -769507651269581073 \\
    0 & 0 & 0 & 0 & 0 & 0 & -97658660653326086254291 & 1539015302539162149 \\
    0 & 0 & 0 & 0 & 0 & 0 & 9403552434308768827094970352 & -148191783481495923038334 \\
  \end{bsmallmatrix}
\]
The SNF of $A$ is given by
\[
  \begin{bsmallmatrix}
    1 & 0 & 0 & 0 & 0 & 0 & 0 & 0 \\
    0 & 1 & 0 & 0 & 0 & 0 & 0 & 0 \\
    0 & 0 & 1 & 0 & 0 & 0 & 0 & 0 \\
    0 & 0 & 0 & 1 & 0 & 0 & 0 & 0 \\
    0 & 0 & 0 & 0 & 1 & 0 & 0 & 0 \\
    0 & 0 & 0 & 0 & 0 & 1 & 0 & 0 \\
    0 & 0 & 0 & 0 & 0 & 0 & 1 & 0 \\
    0 & 0 & 0 & 0 & 0 & 0 & 0 & 10615254 \\
  \end{bsmallmatrix}
\]
The entries of $A$ initially all have modulus at most 10. However intermediate matrices appear with entries of modulus roughly $9.4\times10^{27}$, much larger than any of the entries in the resulting SNF. This ``entry explosion'' is only worsened by increasing the dimensions of the input matrix, and increasing the size of its entries. Arithmetic with entries so large is expensive, so circumventing this issue is of great importance if one wishes to compute the SNF of larger matrices in reasonable time.

\section{Modular techniques}
\label{modular-discussion}

The methods discussed in the previous section to compute Smith Normal Form are inadequate for large matrices. It is not hard to modify the code provided to keep track of the maximum modulus entry.  I constructed a random $100\times100$ matrix whose entries were initially in $\{-1,0,1\}$. After 13 hours of running on the university's servers it has still not terminated. 

The built in {\sc Magma} function \texttt{SmithForm} has no issues computing the Smith Normal Form of such a matrix. This is an indication that the approach taken earlier can be improved. We discuss some ideas described in \cite{sims}, which is heavily inspired by \cite{havas}, which aim to reduce the maximum modulus entry of matrices produced along the way.

We first redevelop some of the theory from the previous sections, modulo a fixed integer. Many of the proofs of the following theorems are similar to ones already given, so we omit them. More detail can be found in \cite{sims}. Fix an integer $d>1$. We use the notation $\overline{A}$ to denote the matrix $A$ with entries reduced modulo $d$.

\begin{definition}
  A \emph{row operation modulo $d$} applied to an integer matrix is one of the following:
  \begin{enumerate}
  \item Swap two rows.
  \item Multiply a row by an integer $c$ with $\gcd(c,d)=1$, reducing entries modulo $d$.
  \item Add an integral multiple of one row to another row, reducing entries modulo $d$.
  \end{enumerate}
  We define \emph{column operations modulo $d$} in a similar fashion. Two $m\times n$ integer matrices $A$ and $B$ are \emph{row equivalent modulo $d$} if there is a sequence of row operations modulo $d$ transforming $\overline{A}$ to $\overline{B}$, and they are \emph{equivalent modulo $d$} if there is a sequence of row and column operations modulo $d$ transforming $\overline{A}$ to $\overline{B}$.
\end{definition}

\begin{definition}
  An integer matrix $A$ is in \emph{row Hermite Normal Form modulo $d$ (HNF modulo $d$)} if
  \begin{enumerate}
  \item $A=\overline{A}$.
  \item $A$ is in row Hermite Normal Form.
  \item The first nonzero entry in each nonzero row divides $d$.
  \end{enumerate}
\end{definition}

Notice that when we work modulo a prime $p$, the third condition ensures that the first nonzero entry in each nonzero row is 1, thus in this case the HNF modulo $p$ agrees with the reduced row echelon form over the field $\Z_p$. In general, the third condition ensures uniqueness in the following theorem.

\begin{theorem}
  If $A$ is an $m\times n$ integer matrix, then there is a unique $m\times n$ integer matrix $B$ with $A\sim B$ and $B$ in row Hermite Normal Form modulo $d$.
\end{theorem}

\begin{definition}
  An integer matrix $A$ is in \emph{Smith Normal Form modulo $d$ (SNF modulo $d$)} if
  \begin{enumerate}
  \item $A=\overline{A}$.
  \item $A$ is in Smith Normal Form.
  \item The nonzero entries of $A$ divide $d$.
  \end{enumerate}
\end{definition}

\begin{theorem}
  If $A$ is an $m\times n$ integer matrix, then there is a unique $m\times n$ integer matrix $B$ with $A\sim B$ and $B$ in Smith Normal Form modulo $d$.
\end{theorem}

The algorithms for computing Hermite and Smith Normal Form modulo an integer are much the same, except all operations are performed modulo $d$, and we multiply rows/columns by appropriate integers $c$ with $\gcd(c,d)=1$ where appropriate to ensure that the relevant entries divide $d$.

\begin{definition}
  Let $A$ be an integer matrix, and let $B$ be the unique matrix in HNF modulo $d$ which is row equivalent modulo $d$ to $A$. The \emph{$d$-rank} of $A$ is the number of nonzero rows of $B$.
\end{definition}

Note that if $A\sim B$ and $B$ is in HNF, then the number of nonzero rows of $B$ is equal to the rank of $A$. Thus the $d$-rank of an integer matrix is a lower bound for its rank.

\begin{definition}
  For an $m\times n$ integer matrix $A$, and $0\leq k\leq m,n$, let $D_k(A)$ denote the gcd of all determinants of $k\times k$ submatrices of $A$. We adopt the convention that the determinant of a $0\times0$ matrix is 1, so that $D_0(A)=1$ for every integer matrix $A$.
\end{definition}

\begin{theorem}
   If integer row or column operations are applied to $A$, then $D_k(A)$ is unchanged for each $0\leq k\leq m,n$.
\end{theorem}

\begin{proof}
  This is clear except when an integer multiple of a row or column is added to another row or column, respectively. See \cite{sims} for the proof.
\end{proof}

\begin{corollary}
  \label{gcdcor}
  Suppose that $A\sim B$, and $B$ is in SNF with invariant factors $d_1\mid d_2\mid\cdots\mid d_r$. Then $D_k(A)=d_1\cdots d_k$ for $k\leq r$, and $D_k(A)=0$ for $k>r$. In particular $d_k=D_k(A)/D_{k-1}(A)$ when $1\leq k\leq r$.
\end{corollary}

This gives an alternate way of computing SNF. Of course, there are $\binom{m}{k}\binom{n}{k}$ different $k\times k$ submatrices of an $m\times n$ matrix, so this quite impractical. However, consider the following.

\begin{theorem}
  \label{recoversnf}
  Let $A$ be an integer matrix, and let $B$ be in SNF with $A\approx B$. Suppose $B$ has invariant factors $d_1\mid d_2\mid \cdots\mid d_r$, and let $d>0$ wtih $d_r\mid d$. Let $C$ be the matrix equivalent to $A$ in SNF modulo $d$. If $C$ has invariant factors $c_1\mid c_2\mid\cdots\mid c_s$, then $s\leq r$, and $d_i=c_i$ for $i\leq s$, and $d_i=d$ for $s<i\leq r$.
\end{theorem}

\begin{proof}
  This follows from noting that $C=\overline{B}$.
\end{proof}

This suggests the following algorithm outline for computing the SNF of a given matrix $A$

\begin{enumerate}
\item Find the rank $r$ of $A$
\item Find a small number of $r\times r$ submatrices of $A$ with nonzero determinant, and take their greatest common divisor $d$. By Corollary \ref{gcdcor}, this will be a multiple of the largest nonzero entry in the SNF of $A$.
\item Compute $C$, the SNF of $A$ modulo $d$
\item If $C$ has $s$ nonzero entries, then by \autoref{recoversnf} we can recover the SNF of $A$ by adding in $r-s$ copies of $d$ on the diagonal after the nonzero entries of $C$.
\end{enumerate}

We do not currently have a method to efficiently compute the rank of a matrix $A$. The HNF algorithm provided earlier also produces relatively large modulus entries for small inputs. We also do not have a method to compute the determinant of an $r\times r$ submatrix of $A$, and find the $r\times r$ submatrices with nonzero determinant.

We can compute the $d$-rank of $A$ by computing the the HNF of $A$ modulo $d$, and this will be a lower bound for the rank of $A$. If we calculate this for various $d$, then their maximum is a lower bound of the rank of $A$. With some work, we can actually ensure that this is equal to the rank.

\begin{theorem}[Hadamard's Inequality]
  Let $A$ be an $n\times n$ real matrix.
  \[ \abs{\det A} \leq \prod_{i=1}^n \left( \sum_{j=1}^n A_{ij}^2 \right)^{1/2} \]
\end{theorem}

\begin{proof}
  See \cite{hadamard}.
\end{proof}

We can use this to get an easily computable bound of the determinant of square submatrix of an $m\times n$ matrix $A$.

\begin{corollary}
  \label{hbound}
  If $A$ is an $m\times n$ integer matrix with $m\leq n$, then the determinant of a square submatrix has modulus at most
  \[ \prod_i \left( \sum_{j=1}^n A_{ij}^2 \right)^{1/2} \text{ and } \left[ \max_j \left( \sum_{i=1}^m A_{ij}^2 \right)^{1/2} \right]^m, \]
  where the product is taken of the nonzero rows.
\end{corollary}

\begin{definition}
  Given an $m\times n$ matrix, we define $h(A)$ to be the bound given by Corollary \ref{hbound} on $A$ if $m\leq n$, and on $A^t$ if $m>n$.
\end{definition}

\begin{theorem}
  Let $A$ be an integer matrix. If $p_1,\ldots,p_k$ are distinct primes with $p_1\cdots p_k>h(A)$, then the rank of $A$ is a the maximum $p_i$-rank of $A$ for $1\leq i\leq k$.
\end{theorem}

\begin{proof}
  Let $r$ be the maximum $p_i$-rank. It suffices to prove that all square submatrices of $A$ with more than $r$ rows have determinant 0. Let $B$ be such a matrix. If $1\leq i\leq k$, since the $p_i$-rank of $A$ is at most $r$, it follows that $\det B\equiv0\pmod{p_i}$. Hence $\det B\equiv0\pmod{p_1\cdots p_k}$. But by Corollary \ref{hbound}, $\abs{\det B}\leq h(A)<p_1\cdots p_k$, thus $\det B=0$.
\end{proof}

Once we have found the rank $r$ of an integer matrix $A$, we can similarly exploit Hadamard's Inequality to efficiently compute the determinants of $r\times r$ submatrices.

\begin{theorem}
  Let $A$ be an integer matrix, and let $B$ be a square submatrix of $A$. Let $p_1,\ldots,p_k$ be distinct primes with $p_1\cdots p_k>2h(A)$. Then if $\det B\equiv b_i\pmod{p_i}$, $1\leq i\leq k$, then $\det B$ is the integer $b$ with smallest absolute value satisfying $b\equiv b_i\pmod{p_i}$, $1\leq i\leq k$.
\end{theorem}

\begin{proof}
  By the Chinese Remainder Theorem, the conditions determine $\det B$ modulo $p_1\cdots p_k$. Also $-h(A)\leq\det B\leq h(A)$, and since $p_1\cdots p_k>2h(A)$, no two integers between $-h(A)$ and $h(A)$ are congruent modulo $p_1\cdots p_k$. 
\end{proof}

Thus we have yet to discuss how to compute the $p$-rank of a matrix, and its determinant modulo $p$, for some prime $p$. However when $p$ is prime, $\Z_p$ is field, so standard methods (e.g. row reduction to row echelon form over a field) can be used to compute the $p$-rank and determinant modulo $p$.

Moreover, \cite{sims} suggests slightly modifying the procedure of row reduction modulo $p$, to find $r\times r$ submatrices of nonzero determinant, where $r$ is the rank. For a matrix $A$ with reduced row echelon form $R$, let $S$ denote the row indices in $A$ which are eventually swapped into a nonzero row position in $R$, and let $T$ denote the column indices which contain the first nonzero entry in some row of $R$. Then the $r\times r$ submatrix of $A$ with rows indexed by $S$ and columns indexed by $T$ will have nonzero determinant.

\section{Empirical Results and Discussion}
\label{conclusion}

Based on the discussions in this report, there are three natural methods available to compute the Smith Normal Form of a given integer matrix in {\sc Magma}. These are
\begin{itemize}
\item \texttt{SmithNormalForm} given in \autoref{snfcode} (described in Section \ref{snf-discussion}), \item \texttt{SmithNormalFormImproved} given in \autoref{snfimprovedcode} (described in Section \ref{modular-discussion}),
\item the {\sc Magma} intrinsic \texttt{SmithForm}.
\end{itemize}

We generated a random $n\times n$ matrix with entries in $\{-1,0,1\}$ for $n=10,20,\ldots,100$, and computed the Smith Normal Form of each of the matrices with each method. The machine used had a clock speed of 2.6GHz, a total of 690GB of RAM, and was running {\sc Magma} version V2.23-3. If a given computation took longer than 30 minutes to halt, it was terminated.

The results of this are recorded in the following table. In each computation, the time taken to compute the SNF, and total memory used in the {\sc Magma} session were recorded. For the computations using \texttt{SmithNormalForm}, the maximum modulus of an entry appearing in some intermediate matrix was recorded (see the ``Max'' column). For the computations using \texttt{SmithNormalFormImproved}, the modulus which calculations were performed with respect to was recorded (see the ``Modulus'' column). There is no corresponding measurement available for the {\sc Magma} intrinsic \texttt{SmithForm}.

  {\centering
  \resizebox{\columnwidth}{!}{%
  \begin{tabular}{|c||c|c|c||c|c|c||c|c|}
    \cline{2-9}
    \multicolumn{1}{c||}{} & \multicolumn{3}{|c||}{\texttt{SmithNormalForm}} & \multicolumn{3}{|c||}{\texttt{SmithNormalFormImproved}} & \multicolumn{2}{|c|}{\texttt{SmithForm}} \\
    \hline
    Dimensions & Time & Memory & Max & Time & Memory & Modulus & Time & Memory \\
    \hline
    $10\times10$ & 0.00s & 32.0MB & 68 & 0.01s & 32.0MB & 40 & 0.00s & 32.0MB \\
    $20\times20$ & 0.00s & 32.0MB & $6.0\cdot10^{42}$ & 0.01s & 32.0MB & $1.7\cdot10^7$ & 0.00s & 32.0MB \\
    $30\times30$ & 0.04s & 32.0MB & $7.8\cdot10^{4056}$ & 0.02s & 32.0MB & $2.5\cdot10^{13}$ & 0.00s & 32.0MB \\
    $40\times40$ & 0.31s & 32.0MB & $\sim10^{1.4\cdot10^5}$ & 0.05s & 32.0MB & $1.3\cdot10^{19}$ & 0.02s & 32.0MB \\
    $50\times50$ & >30m & 384MB & $\sim10^{2.5\cdot10^6}$ & 0.08s & 32.0MB & $3.6\cdot10^{27}$ & 0.02s & 32.0MB \\
    $60\times60$ & >30m & 928MB & $\sim10^{2.5\cdot10^6}$ & 0.15s & 32.0MB & $7.0\cdot10^{34}$ & 0.02s & 32.0MB \\
    $70\times70$ & >30m & 2.89GB & $\sim10^{5.4\cdot10^6}$ & 0.24s & 32.0MB & $5.0\cdot10^{42}$ & 0.02s & 32.0MB \\
    $80\times80$ & >30m & 25.0GB & $\sim10^{1.2\cdot10^7}$ & 0.38s & 32.0MB & $2.0\cdot10^{51}$ & 0.03s & 32.0MB \\
    $90\times90$ & >30m & 2.56GB & $\sim10^{8.9\cdot10^5}$ & 0.60s & 64.0MB & $1.9\cdot10^{60}$ & 0.07s & 32.0MB \\
    $100\times100$ & >30m & 14.8GB & $\sim10^{2.3\cdot10^6}$ & 0.79s & 64.0MB & $1.6\cdot10^{58}$ & 0.06s & 32.0MB \\
    \hline
  \end{tabular}%
  }}\newline


\texttt{SmithNormalFormImproved} outperforms the original procedure \texttt{SmithNormalForm} only when the dimensions of the input are sufficiently large. This is because some extra computation, e.g. the bound given by Corollary \ref{hbound}, is performed by \texttt{SmithNormalFormImproved}. The cost of these additional calculations are only outweighed by the benefit of performing row operations with respect to an appropriate modulus, when the entries appearing in intermediate matrices are very large, which occurs when the input dimensions are large.

Observe that the {\sc Magma} intrinsic \texttt{SmithForm} always performs better than the procedure \texttt{SmithNormalFormImproved}. One reason for this is \texttt{SmithForm} is a compiled C program, however \texttt{SmithNormalFormImproved} is written in {\sc Magma}, an interpreted language.

Moreover, there are other techniques beyond those discussed in \autoref{modular-discussion} which can be used to more efficiently compute the SNF, many of which are discussed in \cite{havas}. Recall the beginning of the proof of \autoref{SNF-existence}, in which a nonzero entry which divided every term in its row and column was needed. If one did not exist, then by choosing a nonzero entry we could produce a smaller entry by applying some row and column operations.

In \cite{havas}, this process is called \emph{pivoting}, and the nonzero entry which is selected is called a \emph{pivot}. In \texttt{SmithNormalForm}, the way a pivot is selected is to initially choose the smallest modulus entry, and thereafter select the first entry in the current pivot's row and column which is not divisible by the current pivot. \cite{havas} suggests some alternate pivoting strategies which can improve performance. In particular it is recommended to try certain pivoting strategies first, and only resort to modular methods if these are unsuccessful.

There is a final deficiency in both of the procedures described in this report. Given an $m\times n$ integer matrix $A$, computing the SNF of $A$ determines the rank $r$ and positive integers $d_1\mid d_2\mid\cdots\mid d_r$ such that $\Z^n/S(A) \cong \Z_{d_1}\oplus\cdots\oplus\Z_{d_r}\oplus\Z^{n-r}$. However the SNF alone does not provide an explicit description of such isomorphism, which is often useful.

When only row and column operations and no modular techniques are used to compute the SNF of $A$, we can find $P\in\GL_m(\Z)$, $Q\in\GL_n(\Z)$ such that $PAQ$ is in SNF, by applying the row operations to $I_m$ to produce $P$, and the column operations to $I_n$ to produce $Q$. As in the proof of \autoref{commutative-diagram}, $P$ and $Q$ can be used to describe such an isomorphism.

Unfortunately, when modular techniques are applied, such an isomorphism can only be recovered in this way if $r=n$. If we define $T\coloneqq\Z_{d_1}\oplus\cdots\oplus\Z_{d_r}$ and $F\coloneqq\Z^{n-r}$, then $\Z^n/S(A)\cong T\oplus F$, and $T$ is torsion (every element has finite order), and $\Z^{n-r}$ is torsion free (all non-identity elements have infinite order). \cite{havas} describes how \emph{Lattice Basis Reduction} algorithms such as {\sc mlll} can be used together with the modular methods described here, to completely describe an isomorphism $\Z^n/S(A)\to T\oplus F$.

\printbibliography
  
\appendix

\section{Hermite Normal Form}

\label{hnfcode}

\lstinputlisting[language=magma,tabsize=8]{../magma/RowReduce.m}

\section{Smith Normal Form}

\label{snfcode}

\lstinputlisting[language=magma,tabsize=8]{../magma/SmithNormalForm.m}

\section{Smith Normal Form Improved}

\label{snfimprovedcode}

\lstinputlisting[language=magma,tabsize=8]{../magma/SNFImproved.m}

\end{document}
